% Luc Destombe --- licence CC BY-NC-SA 4.0
% https://creativecommons.org/licenses/by-nc-sa/4.0/deed.fr
% Fichier autonome : compiler avec pdflatex, deux fois (signets, nombre de pages).
\documentclass[11pt,a4paper]{article}
\makeatletter
\newif\iflycee@unecolonne
\lycee@unecolonnetrue
\newif\iflycee@palatino
\lycee@palatinotrue

\RequirePackage[utf8]{inputenc}
\RequirePackage[T1]{fontenc}
\RequirePackage[french]{babel}
\RequirePackage{etoolbox}

\newcommand{\ShorthandsOff}{\@ifundefined{shorthandoff}{}{\shorthandoff{;:!?}}}
\newcommand{\ShorthandsOn}{\@ifundefined{shorthandon}{}{\shorthandon{;:!?}}}
\ShorthandsOff
\AtBeginDocument{\ShorthandsOn}
\AtBeginEnvironment{tikzpicture}{\ShorthandsOff}

\RequirePackage{geometry}
\RequirePackage{fancyhdr}
\RequirePackage{lastpage}
\RequirePackage{multicol}
\RequirePackage{array}
\RequirePackage{enumitem}
\RequirePackage{xcolor}
\RequirePackage{needspace}

\RequirePackage{amsmath,amssymb}
\RequirePackage{mathpazo}
\DeclareMathSizes{10.95}{10.95}{8}{5.5}
\begingroup\catcode`\_=\active \catcode`\^=\active
\gdef_#1{\sb{\scriptscriptstyle #1}}
\gdef^#1{\sp{\scriptscriptstyle\lycee@moinsepais #1}}
\endgroup
\newcommand\lycee@moins{\mathbin{%
  \setbox\z@\hbox{$\m@th\scriptscriptstyle\lycee@moinsorig$}%
  \dimen@=.55\wd\z@ \advance\dimen@-.5pt
  \hbox to.75\wd\z@{\hss$\m@th\scriptscriptstyle\vcenter{\hbox{%
    \tikz\draw[line width=.5pt,line cap=round](0,0)--(\the\dimen@,0);}}$\hss}}}
\begingroup\lccode`\~=`\- \lowercase{\endgroup\let~}\lycee@moins
\newcommand\lycee@moinsepais{\mathcode`\-="8000 }
\AtBeginDocument{\mathchardef\lycee@moinsorig=\mathcode`\-\relax}
\AtBeginDocument{\catcode`\_=12 \mathcode`\_="8000
  \catcode`\^=12 \mathcode`\^="8000 }
\newcommand\lycee@exposants{%
  \fontdimen13\textfont2=.48\fontdimen6\textfont2
  \fontdimen14\textfont2=.48\fontdimen6\textfont2
  \fontdimen15\textfont2=.48\fontdimen6\textfont2
  \fontdimen13\scriptfont2=.48\fontdimen6\scriptfont2
  \fontdimen14\scriptfont2=.48\fontdimen6\scriptfont2
  \fontdimen15\scriptfont2=.48\fontdimen6\scriptfont2}
\every@math@size\expandafter{\the\every@math@size\lycee@exposants}
\newlength{\retraitcalculs}
\setlength{\retraitcalculs}{2em}

\RequirePackage{tikz}
\usetikzlibrary{arrows.meta,calc,babel,shapes.misc}
\RequirePackage[most]{tcolorbox}
\tcbuselibrary{skins,breakable}

\definecolor{coulDef}{HTML}{1F4E79}
\definecolor{coulProp}{HTML}{7030A0}
\definecolor{coulRem}{HTML}{C55A11}
\definecolor{coulEx}{HTML}{2E7D32}
\definecolor{coulExo}{HTML}{B03A2E}
\definecolor{coulPart}{HTML}{1F4E79}
\definecolor{coulMeth}{HTML}{1B6E4F}
\definecolor{coulCourbe}{HTML}{0B5ED7}
\definecolor{coulCourbeB}{HTML}{C00000}
\definecolor{coulCourbeC}{HTML}{2E7D32}
\definecolor{coulGrille}{HTML}{8C8C8C}
\definecolor{coulRep}{HTML}{1B6E4F}
\definecolor{coulBar}{HTML}{2E7D32}

\newcommand{\R}{\mathbb{R}}
\newcommand{\Rs}{\mathbb{R}^{*}}
\renewcommand{\le}{\leqslant}
\renewcommand{\ge}{\geqslant}
\renewcommand{\approx}{\simeq}
\newcommand{\vect}[1]{\overrightarrow{#1}}
\newcommand{\norme}[1]{\left\lVert #1 \right\rVert}
\newcommand{\intervalle}[2]{\left[#1\,;#2\right]}
\RequirePackage{cancel}
\renewcommand{\CancelColor}{\color{coulExo}}
\newcommand{\fois}[1]{\times{\color{coulExo}#1}}
\newcommand{\barre}[1]{\cancel{#1}}

\tikzset{grille/.style={coulGrille,line width=0.4pt}}
\newcommand{\quadrillage}[4]{\draw[grille,step=1] (#1,#2) grid (#3,#4);}
\tikzset{
  croix/.style={cross out,draw,line width=0.6pt,inner sep=0pt,minimum size=4pt},
  croixdroite/.style={croix,rotate=45,line width=0.8pt},
}
\newcommand{\pt}[3]{\path (#1) node[croix]{} node[#2,font=\small,inner sep=2.5pt]{$#3$};}

\newcommand{\lycee@repere}[4]{%
  \draw[grille,step=1] (#1,#2) grid (#3,#4);
  \draw[-{Stealth[length=2mm]},thick] (#1,0)--({#3+0.4},0) node[below left=-1pt]{$x$};
  \draw[-{Stealth[length=2mm]},thick] (0,#2)--(0,{#4+0.4}) node[below left=-1pt]{$y$};
  \node[below left,font=\scriptsize,inner sep=1.5pt] at (0,0) {$0$};}
\newcommand{\lycee@gradx}[1]{\foreach \k/\l in {#1}{%
  \draw (\k,0.07)--(\k,-0.07) node[below,font=\scriptsize,inner sep=1.5pt]{$\l$};}}
\newcommand{\lycee@grady}[1]{\foreach \k/\l in {#1}{%
  \draw (0.07,\k)--(-0.07,\k) node[left,font=\scriptsize,inner sep=1.5pt]{$\l$};}}
\newcommand{\lycee@intff}[2]{\left[#1\,;#2\right]}
\newcommand{\lycee@intfo}[2]{\left[#1\,;#2\right[}
\newcommand{\lycee@intof}[2]{\left]#1\,;#2\right]}
\newcommand{\lycee@intoo}[2]{\left]#1\,;#2\right[}
\newcommand{\lycee@ens}[1]{\left\{#1\right\}}
\AtBeginDocument{%
  \providecommand{\repere}{\lycee@repere}%
  \providecommand{\gradx}{\lycee@gradx}%
  \providecommand{\grady}{\lycee@grady}%
  \providecommand{\Cf}{\mathcal{C}\sb{\scriptscriptstyle f}}%
  \providecommand{\Cg}{\mathcal{C}\sb{\scriptscriptstyle g}}%
  \providecommand{\intff}{\lycee@intff}%
  \providecommand{\intfo}{\lycee@intfo}%
  \providecommand{\intof}{\lycee@intof}%
  \providecommand{\intoo}{\lycee@intoo}%
  \providecommand{\ens}{\lycee@ens}}

\newlist{questions}{enumerate}{3}
\setlist[questions,1]{label=\arabic*\textdegree),leftmargin=*,itemsep=2pt,topsep=3pt}
\setlist[questions,2]{label=\alph*),leftmargin=*,itemsep=1pt,topsep=2pt}
\setlist[questions,3]{label=\roman*),leftmargin=*,itemsep=1pt,topsep=1pt}
\setlist[itemize]{label=\textbullet,leftmargin=*,itemsep=2pt,topsep=3pt}

\newcommand{\besoinplace}[1]{\par\penalty\@M\begingroup
  \dimen@=#1\relax\dimen@ii=\pagegoal\advance\dimen@ii-\pagetotal
  \ifdim\dimen@>\dimen@ii\ifdim\dimen@ii>\z@\vfil\fi\break\fi\endgroup}

\newcommand{\lycee@pied}{\thepage/\pageref{LastPage}}
\newcommand{\entete}[2]{%
  \pagestyle{fancy}\fancyhf{}%
  \fancyhead[L]{\footnotesize\color{coulPart}\textbf{#1}}%
  \fancyhead[R]{\footnotesize\color{coulPart}\textbf{#2}}%
  \fancyfoot[C]{\footnotesize\lycee@pied}%
  \renewcommand{\headrulewidth}{0.4pt}%
  \renewcommand{\headrule}{\hbox to\headwidth{%
    \color{coulPart!40}\leaders\hrule height \headrulewidth\hfill}}}

\RequirePackage{ccicons}
\newcommand{\lycee@licence}{{\scriptsize\color{gray}%
  \copyright\ Luc Destombe --- \ccbyncsa\ CC BY-NC-SA 4.0}}
\newif\iflycee@licence \lycee@licencetrue
\newcommand{\sanslicence}{\lycee@licencefalse}
\AtBeginDocument{\iflycee@licence\AddToHookNext{shipout/foreground}{%
  \put(0,0){\rlap{\kern\dimexpr1in+\hoffset+\oddsidemargin+\textwidth\relax
    \llap{\raisebox{-\dimexpr1in+\voffset+\topmargin+\headheight+\headsep
      +\textheight+\footskip\relax}{\lycee@licence}}}}}\fi}

\newcommand{\formulesserrees}{%
  \setlength{\abovedisplayskip}{3pt}\setlength{\belowdisplayskip}{3pt}%
  \setlength{\abovedisplayshortskip}{2pt}\setlength{\belowdisplayshortskip}{3pt}}

\setlength{\parindent}{0pt}

\RequirePackage[implicit=false,hidelinks,pdfencoding=auto,psdextra,bookmarksopen,
  bookmarksopenlevel=1,pdfcreator={},pdfproducer={}]{hyperref}
\RequirePackage{bookmark}
\hypersetup{pdfauthor={Luc Destombe},pdfkeywords={CC BY-NC-SA 4.0}}
\newcounter{lycee@signet}
\newcommand{\signet}[3][]{\stepcounter{lycee@signet}%
  \edef\lycee@dest{\if\relax\detokenize{#1}\relax
    signet.\arabic{lycee@signet}\else#1\fi}%
  \hypertarget{\lycee@dest}{}\bookmark[level=#2,dest=\lycee@dest]{#3}}
\pdfstringdefDisableCommands{%
  \def\R{R}\def\vect#1{#1}\def\Cf{Cf}\def\Cg{Cg}\def\fois#1{×#1}%
  \def\barre#1{#1}\def\intervalle#1#2{[#1;#2]}\def\intff#1#2{[#1;#2]}%
  \def\textsuperscript#1{#1}\def\dfrac#1#2{#1/#2}\def\frac#1#2{#1/#2}%
  \def\vec#1{#1⃗}}%
\RequirePackage[official]{eurosym}

\tcbset{exercice corrige/.style={
  enhanced,breakable,before upper=\raggedright,
  colback=white,colframe=coulExo,
  boxrule=0.8pt,arc=2pt,
  left=6pt,right=6pt,top=8pt,bottom=5pt,
  before skip=9pt,after skip=4pt,
  coltitle=white,fonttitle=\bfseries\small,
  attach boxed title to top left={xshift=8pt,yshift=-\tcboxedtitleheight/2},
  boxed title style={colback=coulExo,arc=2pt,boxrule=0pt}}}
\newcounter{exo}
\newtcolorbox{exercice}[1][]{exercice corrige,
  phantom={\signet[exercice-\arabic{exo}]{2}{Exercice \arabic{exo}}},
  title={Exercice \theexo\ifx\relax#1\relax\else{} --- #1\fi}}
\newenvironment{exo}[1][\relax]{\refstepcounter{exo}\begin{exercice}[#1]}{\end{exercice}}

\geometry{top=1.8cm,bottom=1cm,left=1cm,right=1cm,headheight=14pt,footskip=0.6cm}
\renewcommand{\lycee@pied}{\thepage}

\renewcommand{\theexo}{\ifnum\value{exo}<10 0\fi\arabic{exo}}
\tcbset{exercice corrige/.style={
  enhanced,breakable,before upper=\raggedright,
  colback=white,colframe=coulExo,
  boxrule=0.9pt,arc=2pt,
  left=8pt,right=8pt,top=8pt,bottom=6pt,
  before skip=8pt,after skip=6pt,
  coltitle=white,fonttitle=\bfseries,
  attach boxed title to top left={xshift=8pt,yshift=-\tcboxedtitleheight/2},
  boxed title style={colback=coulExo,arc=2pt,boxrule=0pt}}}

\newcommand{\entetefiche}[2]{%
  \begin{center}
    \begin{tcolorbox}[enhanced,width=0.92\linewidth,colback=white,colframe=coulPart,
        boxrule=1.4pt,arc=3pt,halign=center,top=6pt,bottom=6pt]
      {\large\bfseries\color{coulPart} CORRIGÉ --- FICHE D'EXERCICES}\\[3pt]
      {\normalsize\bfseries #1}\\[2pt]
      {\normalsize #2}
    \end{tcolorbox}
  \end{center}}

\newcounter{partie}
\newcommand{\partie}[1]{%
  \stepcounter{partie}%
  \besoinplace{8\baselineskip}\vspace{4pt}%
  \begin{tcolorbox}[enhanced,colback=coulPart!8,colframe=coulPart,boxrule=0pt,
      leftrule=3pt,arc=0pt,left=6pt,right=4pt,top=3pt,bottom=3pt,
      before skip=6pt,after skip=2pt]
    \bfseries\color{coulPart}\Roman{partie} --- #1\signet{1}{\Roman{partie} --- #1}
  \end{tcolorbox}}
\newcommand{\partiesansnum}[1]{%
  \besoinplace{8\baselineskip}\vspace{4pt}%
  \begin{tcolorbox}[enhanced,colback=coulPart!8,colframe=coulPart,boxrule=0pt,
      leftrule=3pt,arc=0pt,left=6pt,right=4pt,top=3pt,bottom=3pt,
      before skip=6pt,after skip=2pt]
    \bfseries\color{coulPart}#1\signet{1}{#1}
  \end{tcolorbox}}

\setlist[questions,1]{label=\arabic*\textdegree),leftmargin=*,itemsep=3pt,topsep=3pt}
\setlist[questions,2]{label=\alph*),leftmargin=*,itemsep=2pt,topsep=2pt}
\setlist[questions,3]{label=\roman*),leftmargin=*,itemsep=1pt,topsep=1pt}
\newlist{expr}{enumerate}{1}
\setlist[expr]{label={\Alph*\ $=$},leftmargin=*,itemsep=2pt,topsep=3pt}

\newcommand{\res}[1]{{\color{coulRep}#1}}
\newcommand{\rem}[1]{\par\smallskip{\small\itshape\color{black!60}$\rightarrow$ #1}}
\newenvironment{calculs}{\par\smallskip\noindent$\displaystyle\begin{aligned}[t]}%
  {\end{aligned}$\par\smallskip}
\newcommand{\justif}[1]{\quad\text{\small\itshape\color{black!60}#1}}
\newcommand{\cas}[1]{\par\smallskip\noindent\textbf{\color{coulExo}#1}\nobreak}
\newcommand{\explique}[2]{\par\smallskip\noindent
  \begin{minipage}[t]{0.6\linewidth}\raggedright #1\end{minipage}\hfill
  \begin{minipage}[t]{0.37\linewidth}\raggedright\leavevmode
    {\small\itshape\color{black!60}#2}\end{minipage}\par}

\setlength{\parskip}{2pt}

\makeatother
\geometry{top=1.8cm,bottom=1.6cm,left=1.8cm,right=1.8cm,headheight=14pt,footskip=30pt}
\entete{Chapitre 1 --- Calcul littéral --- corrigé}{Première spécialité}
\setlength{\parindent}{0pt}
\setlength{\jot}{5pt}

\let\casligne\cas
\renewcommand{\cas}[1]{\casligne{#1}\par\nopagebreak}

\usepackage{tkz-tab}

\begin{document}

\entetefiche{Chapitre 1 --- Calcul littéral}{Première spécialité mathématiques}

\partie{Réduire, développer, identités remarquables}

\begin{exo}[Réduire et distribuer]
\explique{Deux gestes seulement : \emph{réduire} une somme (regrouper les termes
semblables) et \emph{distribuer} un facteur sur une parenthèse, $k(a+b)=ka+kb$.}{méthodes 1
et 2 du cours}
\begin{calculs}
A &= 3x+8+6x &&\justif{on recopie}\\
A &= 3x+6x+8 &&\justif{on rapproche les termes semblables : les termes en $x$}\\
A &= \res{9x+8} &&\justif{$3+6=9$ ; $9x$ et $8$ ne sont pas semblables : on s'arrête}
\end{calculs}
\begin{calculs}
B &= 4x^{2}-9x^{2}+2x^{2} &&\justif{trois termes en $x^{2}$ : tous semblables}\\
B &= \res{-3x^{2}} &&\justif{on additionne les coefficients : $4-9+2=-3$}
\end{calculs}
\begin{calculs}
C &= 6(3x-2) &&\justif{on recopie}\\
C &= 6\times 3x-6\times 2 &&\justif{le $6$ multiplie \emph{chaque} terme de la parenthèse}\\
C &= \res{18x-12} &&
\end{calculs}
\begin{calculs}
D &= 2x(5x-4) &&\justif{on recopie}\\
D &= 2x\times 5x-2x\times 4 &&\justif{$2x$ multiplie chaque terme}\\
D &= \res{10x^{2}-8x} &&\justif{$2\times 5=10$ et $x\times x=x^{2}$}
\end{calculs}
\begin{calculs}
E &= 9(x-1)-4x &&\justif{on distribue d'abord le $9$ : le produit passe avant la
soustraction}\\
E &= 9x-9-4x &&\\
E &= 9x-4x-9 &&\justif{on rapproche les termes en $x$}\\
E &= \res{5x-9} &&
\end{calculs}
\begin{calculs}
F &= -5x+3(2x+7) &&\justif{on distribue le $3$}\\
F &= -5x+6x+21 &&\\
F &= \res{x+21} &&\justif{$-5+6=1$, et $1x$ s'écrit $x$}
\end{calculs}
\begin{calculs}
G &= -(4-3x) &&\justif{le \og $-$\fg{} devant la parenthèse change le signe de
\emph{chacun} de ses termes}\\
G &= -4+3x &&\justif{$-(4)$ donne $-4$ ; $-(-3x)$ donne $+3x$}\\
G &= \res{3x-4} &&\justif{on ordonne : le terme en $x$ d'abord}
\end{calculs}
\begin{calculs}
H &= -2x(1-4x) &&\justif{on distribue $-2x$, \emph{avec} son signe}\\
H &= -2x\times 1+(-2x)\times(-4x) &&\\
H &= -2x+8x^{2} &&\justif{$(-2x)\times(-4x)$ est positif : deux signes $-$}\\
H &= \res{8x^{2}-2x} &&\justif{on ordonne suivant les puissances décroissantes de $x$}
\end{calculs}
\rem{Erreurs fréquentes : en G, écrire $-4-3x$ (on n'a changé que le premier signe) ; en
H, oublier que $(-2x)\times(-4x)$ est positif. En A, $9x+8$ ne devient pas $17x$ : on
n'ajoute que des termes semblables.}
\end{exo}

\begin{exo}[Double distributivité et identités remarquables]
\explique{Avant de développer, on regarde la forme : deux parenthèses différentes, c'est
la double distributivité ; un carré $(a\pm b)^{2}$ ou un produit $(a+b)(a-b)$, c'est une
identité remarquable, plus rapide.}{propriété 1 du cours}
\begin{calculs}
A &= (2x+3)(4x+1) &&\justif{deux parenthèses différentes : pas d'identité remarquable}\\
A &= 2x\times 4x+2x\times 1+3\times 4x+3\times 1 &&\justif{chaque terme de la première
par chaque terme de la seconde : $4$ produits}\\
A &= 8x^{2}+2x+12x+3 &&\\
A &= \res{8x^{2}+14x+3} &&\justif{on réduit : $2x+12x=14x$}
\end{calculs}
\begin{calculs}
B &= (3x-4)(-2x+5) &&\justif{chaque terme garde son signe}\\
B &= 3x\times(-2x)+3x\times 5+(-4)\times(-2x)+(-4)\times 5 &&\\
B &= -6x^{2}+15x+8x-20 &&\justif{$(-4)\times(-2x)=+8x$}\\
B &= \res{-6x^{2}+23x-20} &&\justif{on réduit}
\end{calculs}
\begin{calculs}
C &= (5x+2)^{2} &&\justif{identité remarquable $(a+b)^{2}$, avec $a=5x$ et $b=2$}\\
C &= (5x)^{2}+2\times 5x\times 2+2^{2} &&\justif{$a^{2}+2ab+b^{2}$}\\
C &= \res{25x^{2}+20x+4} &&\justif{$(5x)^{2}=25x^{2}$ : le carré porte aussi sur le $5$}
\end{calculs}
\begin{calculs}
D &= (4-3x)^{2} &&\justif{identité remarquable $(a-b)^{2}$, avec $a=4$ et $b=3x$}\\
D &= 4^{2}-2\times 4\times 3x+(3x)^{2} &&\justif{$a^{2}-2ab+b^{2}$}\\
D &= 16-24x+9x^{2} &&\\
D &= \res{9x^{2}-24x+16} &&\justif{on ordonne}
\end{calculs}
\begin{calculs}
E &= (3x+5)(3x-5) &&\justif{identité remarquable $(a+b)(a-b)$, avec $a=3x$ et $b=5$}\\
E &= (3x)^{2}-5^{2} &&\justif{$a^{2}-b^{2}$}\\
E &= \res{9x^{2}-25} &&\justif{pas de terme en $x$ : les doubles produits se
compensent}
\end{calculs}
\begin{calculs}
F &= \left(x-\dfrac{1}{3}\right)^{2} &&\justif{identité remarquable $(a-b)^{2}$, avec
$a=x$ et $b=\frac{1}{3}$}\\
F &= x^{2}-2\times x\times\dfrac{1}{3}+\left(\dfrac{1}{3}\right)^{2} &&\justif{$a^{2}-2ab+b^{2}$}\\
F &= \res{x^{2}-\dfrac{2}{3}x+\dfrac{1}{9}} &&\justif{$2\times\frac{1}{3}=\frac{2}{3}$ et
$\left(\frac{1}{3}\right)^{2}=\frac{1}{9}$}
\end{calculs}
\begin{calculs}
G &= 3(x+2)(1-x) &&\justif{trois facteurs : on développe d'abord les deux parenthèses}\\
G &= 3\left(x-x^{2}+2-2x\right) &&\justif{double distributivité, dans la parenthèse}\\
G &= 3\left(-x^{2}-x+2\right) &&\justif{on réduit dans la parenthèse}\\
G &= \res{-3x^{2}-3x+6} &&\justif{on distribue le $3$ en dernier}
\end{calculs}
\rem{Erreur fréquente : $(5x+2)^{2}=25x^{2}+4$. Le double produit $2ab$ ne disparaît
pas ; seul $(a+b)(a-b)$ n'a pas de terme en $x$. Contrôle possible : avec $x=1$,
$(5+2)^{2}=49$ et $25+20+4=49$.}
\end{exo}

\begin{exo}[Développer, réduire, ordonner]
\explique{Même démarche pour les cinq expressions : on développe chaque produit à part,
on le garde \emph{entre parenthèses} s'il suit un \og $-$\fg, puis on supprime les
parenthèses en changeant les signes, et on réduit.}{exemple c) de la méthode 2}
\cas{A}
\begin{calculs}
A &= (4x+1)(2x+3)-(5x-2)(x+4) &&\justif{deux doubles distributivités}\\
A &= \left(8x^{2}+12x+2x+3\right)-\left(5x^{2}+20x-2x-8\right) &&\justif{chaque produit
entre parenthèses}\\
A &= \left(8x^{2}+14x+3\right)-\left(5x^{2}+18x-8\right) &&\justif{on réduit dans chaque
parenthèse}\\
A &= 8x^{2}+14x+3-5x^{2}-18x+8 &&\justif{le \og $-$\fg{} change les signes des trois
termes}\\
A &= \res{3x^{2}-4x+11} &&\justif{on regroupe : les $x^{2}$, les $x$, les nombres}
\end{calculs}
\cas{B}
\begin{calculs}
B &= (2x-5)^{2}-(3x+1)(x-2) &&\justif{une identité remarquable, puis un produit}\\
B &= \left(4x^{2}-20x+25\right)-\left(3x^{2}-6x+x-2\right) &&\justif{$(a-b)^{2}$ avec
$a=2x$ et $b=5$}\\
B &= \left(4x^{2}-20x+25\right)-\left(3x^{2}-5x-2\right) &&\\
B &= 4x^{2}-20x+25-3x^{2}+5x+2 &&\justif{le \og $-$\fg{} change les trois signes}\\
B &= \res{x^{2}-15x+27} &&
\end{calculs}
\cas{C}
\explique{Calcul long : on développe d'abord chaque produit à part.}{un bloc par
produit}
\begin{calculs}
\left(x^{2}-x+2\right)(3x+1) &= 3x^{3}+x^{2}-3x^{2}-x+6x+2 &&\justif{chaque terme du
trinôme par $3x$, puis par $1$ : $6$ produits}\\
\left(x^{2}-x+2\right)(3x+1) &= 3x^{3}-2x^{2}+5x+2 &&\justif{on réduit}\\[4pt]
(x+1)(x-1)^{2} &= (x+1)(x-1)(x-1) &&\justif{le carré, c'est $(x-1)$ deux fois}\\
(x+1)(x-1)^{2} &= \left(x^{2}-1\right)(x-1) &&\justif{$(x+1)(x-1)=x^{2}-1$ : identité
remarquable}\\
(x+1)(x-1)^{2} &= x^{3}-x^{2}-x+1 &&
\end{calculs}
\explique{Puis on soustrait le second résultat du premier.}{}
\begin{calculs}
C &= \left(3x^{3}-2x^{2}+5x+2\right)-\left(x^{3}-x^{2}-x+1\right) &&\justif{on garde
les parenthèses}\\
C &= 3x^{3}-2x^{2}+5x+2-x^{3}+x^{2}+x-1 &&\justif{le \og $-$\fg{} change les quatre
signes}\\
C &= \res{2x^{3}-x^{2}+6x+1} &&\justif{ordonné : $x^{3}$, puis $x^{2}$, puis $x$}
\end{calculs}
\cas{D}
\begin{calculs}
D &= (x+2)\left(x^{2}-2x+4\right) &&\justif{$x$ par chaque terme, puis $2$ par chaque
terme}\\
D &= x^{3}-2x^{2}+4x+2x^{2}-4x+8 &&\\
D &= \res{x^{3}+8} &&\justif{les termes en $x^{2}$ et en $x$ s'annulent}
\end{calculs}
\cas{E}
\begin{calculs}
E &= 3x(2-x)-(x-4)(2x+3) &&\\
E &= \left(6x-3x^{2}\right)-\left(2x^{2}+3x-8x-12\right) &&\justif{chaque produit
entre parenthèses}\\
E &= \left(6x-3x^{2}\right)-\left(2x^{2}-5x-12\right) &&\\
E &= 6x-3x^{2}-2x^{2}+5x+12 &&\justif{le \og $-$\fg{} change les trois signes}\\
E &= \res{-5x^{2}+11x+12} &&\justif{on ordonne : $x^{2}$ d'abord}
\end{calculs}
\rem{Le piège commun : le \og $-$\fg{} devant le second produit. Erreur fréquente en E :
écrire $-2x^{2}-5x-12$, en ne changeant que le premier signe. On garde les parenthèses
jusqu'au bout du développement, et on ne les supprime qu'à la ligne suivante.}
\end{exo}

\partie{Substituer une valeur}

\begin{exo}[Calculer une image]
\cas{1) a)}
\begin{calculs}
P(1) &= 3\times 1^{2}-4\times 1-2 &&\justif{on remplace chaque $x$ par $1$}\\
P(1) &= 3-4-2 &&\\
P(1) &= \res{-3} &&\\[4pt]
P(-1) &= 3\times(-1)^{2}-4\times(-1)-2 &&\justif{on remplace $x$ par $-1$,
\emph{entre parenthèses}}\\
P(-1) &= 3\times 1+4-2 &&\justif{$(-1)^{2}=1$ et $-4\times(-1)=+4$}\\
P(-1) &= \res{5} &&
\end{calculs}
\cas{1) b)}
\begin{calculs}
P\!\left(\dfrac{1}{2}\right) &= 3\times\left(\dfrac{1}{2}\right)^{2}-4\times\dfrac{1}{2}-2
  &&\justif{on remplace $x$ par $\frac{1}{2}$}\\
P\!\left(\dfrac{1}{2}\right) &= 3\times\dfrac{1}{4}-2-2 &&\justif{le carré d'abord :
$\left(\frac{1}{2}\right)^{2}=\frac{1}{4}$}\\
P\!\left(\dfrac{1}{2}\right) &= \dfrac{3}{4}-\dfrac{16}{4} &&\justif{$-2-2=-4$, et
$4=\frac{16}{4}$ : même dénominateur}\\
P\!\left(\dfrac{1}{2}\right) &= \res{-\dfrac{13}{4}} &&\\[4pt]
P\!\left(-\dfrac{2}{3}\right) &= 3\times\left(-\dfrac{2}{3}\right)^{2}-4\times
  \left(-\dfrac{2}{3}\right)-2 &&\justif{$-\frac{2}{3}$ entre parenthèses}\\
P\!\left(-\dfrac{2}{3}\right) &= 3\times\dfrac{4}{9}+\dfrac{8}{3}-2
  &&\justif{un carré est positif : $\left(-\frac{2}{3}\right)^{2}=\frac{4}{9}$}\\
P\!\left(-\dfrac{2}{3}\right) &= \dfrac{4}{3}+\dfrac{8}{3}-2
  &&\justif{$3\times\frac{4}{9}=\frac{12}{9}$, qui se simplifie par $3$}\\
P\!\left(-\dfrac{2}{3}\right) &= \dfrac{12}{3}-2 &&\\
P\!\left(-\dfrac{2}{3}\right) &= \res{2} &&\justif{$\frac{12}{3}=4$}
\end{calculs}
\rem{Erreur fréquente : calculer $3\times\left(-\frac{2}{3}\right)^{2}$ comme
$\left(3\times\left(-\frac{2}{3}\right)\right)^{2}$. Le carré passe avant le produit :
il ne porte que sur $-\frac{2}{3}$.}
\cas{2) a)}
\explique{On remplace $x$ par $0$ dans chacune des deux expressions.}{on calcule $A$ et
$B$ à part}
\begin{calculs}
A &= (0+2)(0-5) &&\\
A &= 2\times(-5) &&\\
A &= \res{-10} &&\\[4pt]
B &= 0^{2}+3\times 0-10 &&\\
B &= \res{-10} &&
\end{calculs}
\cas{2) b)}
\begin{calculs}
A &= (1+2)(1-5) &&\justif{on remplace $x$ par $1$}\\
A &= 3\times(-4) &&\\
A &= \res{-12} &&\\[4pt]
B &= 1^{2}+3\times 1-10 &&\\
B &= 1+3-10 &&\\
B &= \res{-6} &&
\end{calculs}
\cas{2) c)}
\explique{\res{Non} : pour $x=1$, on trouve $A=-12$ et $B=-6$. Les deux expressions n'ont
pas toujours la même valeur, donc on ne peut pas affirmer que $A=B$ pour tout réel
$x$.}{un seul contre-exemple suffit pour répondre non}
\rem{Pour $x=0$, les deux valeurs étaient égales : une valeur commune ne prouve rien. En
développant, $A=x^{2}-5x+2x-10$, soit $A=x^{2}-3x-10$ : $A$ et $B$ diffèrent par le
terme en $x$ ($-3x$ au lieu de $3x$).}
\end{exo}

\begin{exo}[Substituer une racine carrée]
\explique{Un seul réflexe nouveau : $\left(\sqrt{a}\right)^{2}=a$. Pour le reste, on
calcule comme d'habitude, puis on regroupe d'un côté les nombres sans racine, de l'autre
les termes en $\sqrt{a}$.}{les racines carrées se développent comme les autres
nombres}
\cas{a)}
\begin{calculs}
P\!\left(\sqrt{5}\right) &= \left(\sqrt{5}\right)^{2}-5 &&\justif{on remplace $x$ par
$\sqrt{5}$}\\
P\!\left(\sqrt{5}\right) &= 5-5 &&\justif{$\left(\sqrt{5}\right)^{2}=5$}\\
P\!\left(\sqrt{5}\right) &= \res{0} &&
\end{calculs}
\cas{b)}
\begin{calculs}
P\!\left(\sqrt{2}\right) &= 2\left(\sqrt{2}\right)^{2}+3\sqrt{2}-1 &&\justif{le carré ne
porte que sur $\sqrt{2}$}\\
P\!\left(\sqrt{2}\right) &= 2\times 2+3\sqrt{2}-1 &&\justif{$\left(\sqrt{2}\right)^{2}=2$}\\
P\!\left(\sqrt{2}\right) &= \res{3+3\sqrt{2}} &&\justif{$4-1=3$ ; $3$ et $3\sqrt{2}$ ne
se regroupent pas}
\end{calculs}
\cas{c)}
\begin{calculs}
P\!\left(\sqrt{3}\right) &= \left(3\sqrt{3}+1\right)\left(\sqrt{3}-2\right) &&\justif{on
remplace dans chaque facteur}\\
P\!\left(\sqrt{3}\right) &= 3\sqrt{3}\times\sqrt{3}-3\sqrt{3}\times 2+\sqrt{3}-2
  &&\justif{double distributivité}\\
P\!\left(\sqrt{3}\right) &= 9-6\sqrt{3}+\sqrt{3}-2 &&\justif{$3\sqrt{3}\times\sqrt{3}=3\times 3$}\\
P\!\left(\sqrt{3}\right) &= \res{7-5\sqrt{3}} &&\justif{$9-2=7$ et
$-6\sqrt{3}+\sqrt{3}=-5\sqrt{3}$}
\end{calculs}
\cas{d)}
\begin{calculs}
P\!\left(\sqrt{5}\right) &= \left(2\sqrt{5}-1\right)^{2}+4\sqrt{5} &&\justif{$(a-b)^{2}$,
avec $a=2\sqrt{5}$ et $b=1$}\\
P\!\left(\sqrt{5}\right) &= \left(2\sqrt{5}\right)^{2}-2\times 2\sqrt{5}\times 1+1^{2}+4\sqrt{5}
  &&\justif{$a^{2}-2ab+b^{2}$}\\
P\!\left(\sqrt{5}\right) &= 20-4\sqrt{5}+1+4\sqrt{5} &&\justif{$\left(2\sqrt{5}\right)^{2}=4\times 5$}\\
P\!\left(\sqrt{5}\right) &= \res{21} &&\justif{les termes en $\sqrt{5}$ s'annulent}
\end{calculs}
\cas{e)}
\begin{calculs}
P\!\left(1+\sqrt{2}\right) &= \left(1+\sqrt{2}\right)^{2}-2\left(1+\sqrt{2}\right)+3
  &&\justif{$x$ remplacé par $1+\sqrt{2}$, entre parenthèses}\\
P\!\left(1+\sqrt{2}\right) &= 1+2\sqrt{2}+2-2-2\sqrt{2}+3 &&\justif{$(a+b)^{2}$ :
$1+2\sqrt{2}+\left(\sqrt{2}\right)^{2}$ ; on distribue $-2$}\\
P\!\left(1+\sqrt{2}\right) &= \res{4} &&\justif{$1+2-2+3=4$ ; les $\sqrt{2}$
s'annulent}
\end{calculs}
\rem{Erreur fréquente : $\left(2\sqrt{5}\right)^{2}=2\times 5$. Le carré porte sur les deux
facteurs : $2^{2}\times\left(\sqrt{5}\right)^{2}=4\times 5$.}
\end{exo}

\begin{exo}[Racines carrées dans un quotient]
\cas{1) a)}
\explique{$x-3=0$ pour $x=3$ : la valeur interdite de $f$ est \res{$3$}.}{un dénominateur
ne peut pas être nul : on ne peut pas diviser par $0$}
\cas{1) b)}
\begin{calculs}
f(5) &= \dfrac{5+1}{5-3} &&\justif{numérateur et dénominateur calculés à part}\\
f(5) &= \dfrac{6}{2} &&\\
f(5) &= \res{3} &&\\[4pt]
f\!\left(\dfrac{1}{2}\right) &= \dfrac{\frac{1}{2}+1}{\frac{1}{2}-3} &&\justif{$1=\frac{2}{2}$
et $3=\frac{6}{2}$}\\
f\!\left(\dfrac{1}{2}\right) &= \dfrac{\frac{3}{2}}{-\frac{5}{2}} &&\\
f\!\left(\dfrac{1}{2}\right) &= \dfrac{3}{\barre{2}}\times\left(-\dfrac{\barre{2}}{5}\right)
  &&\justif{diviser par $-\frac{5}{2}$, c'est multiplier par $-\frac{2}{5}$}\\
f\!\left(\dfrac{1}{2}\right) &= \res{-\dfrac{3}{5}} &&
\end{calculs}
\cas{1) c)}
\explique{Le dénominateur $\sqrt{5}-3$ contient une racine : on multiplie en haut et en bas
par sa \emph{quantité conjuguée} $\sqrt{5}+3$ (même expression, signe du milieu
changé).}{avec $(a-b)(a+b)=a^{2}-b^{2}$, la racine est élevée au carré et disparaît}
\begin{calculs}
f\!\left(\sqrt{5}\right) &= \dfrac{\sqrt{5}+1}{\sqrt{5}-3} &&\justif{il reste une racine
au dénominateur}\\
f\!\left(\sqrt{5}\right) &= \dfrac{\left(\sqrt{5}+1\right)\left(\sqrt{5}+3\right)}
  {\left(\sqrt{5}-3\right)\left(\sqrt{5}+3\right)} &&\justif{on multiplie en haut et en
bas par $\sqrt{5}+3$}\\
\color{black!60}\left(\sqrt{5}-3\right)\left(\sqrt{5}+3\right) &\color{black!60}= 5-9
  &&\justif{calcul intermédiaire : $a^{2}-b^{2}$}\\
\color{black!60}\left(\sqrt{5}+1\right)\left(\sqrt{5}+3\right) &\color{black!60}=
  5+3\sqrt{5}+\sqrt{5}+3 &&\justif{calcul intermédiaire : double distributivité}\\
f\!\left(\sqrt{5}\right) &= \dfrac{8+4\sqrt{5}}{-4} &&\\
f\!\left(\sqrt{5}\right) &= \res{-2-\sqrt{5}} &&\justif{on divise chaque terme du
numérateur par $-4$}
\end{calculs}
\cas{2)}
\explique{Le dénominateur $\sqrt{3}+1$ contient une racine : on multiplie par sa quantité
conjuguée $\sqrt{3}-1$.}{même méthode qu'en 1) c)}
\begin{calculs}
g\!\left(\sqrt{3}\right) &= \dfrac{\sqrt{3}-2}{\sqrt{3}+1} &&\justif{on remplace $x$ par
$\sqrt{3}$}\\
g\!\left(\sqrt{3}\right) &= \dfrac{\left(\sqrt{3}-2\right)\left(\sqrt{3}-1\right)}
  {\left(\sqrt{3}+1\right)\left(\sqrt{3}-1\right)} &&\justif{quantité conjuguée en haut et
en bas}\\
\color{black!60}\left(\sqrt{3}+1\right)\left(\sqrt{3}-1\right) &\color{black!60}= 3-1
  &&\justif{calcul intermédiaire : $a^{2}-b^{2}$}\\
\color{black!60}\left(\sqrt{3}-2\right)\left(\sqrt{3}-1\right) &\color{black!60}=
  3-\sqrt{3}-2\sqrt{3}+2 &&\justif{calcul intermédiaire}\\
g\!\left(\sqrt{3}\right) &= \res{\dfrac{5-3\sqrt{3}}{2}} &&\justif{plus de racine au
dénominateur}
\end{calculs}
\cas{3)}
\explique{On remplace $x$ par $2+\sqrt{3}$ : le dénominateur devient
$\left(2+\sqrt{3}\right)-2$, une racine seule. On multiplie alors en haut et en bas par
cette racine.}{$\sqrt{3}\times\sqrt{3}=3$ : la racine disparaît}
\begin{calculs}
h\!\left(2+\sqrt{3}\right) &= \dfrac{2+\sqrt{3}+3}{2+\sqrt{3}-2} &&\justif{on remplace
en haut et en bas}\\
h\!\left(2+\sqrt{3}\right) &= \dfrac{5+\sqrt{3}}{\sqrt{3}} &&\justif{$2-2=0$ au
dénominateur}\\
h\!\left(2+\sqrt{3}\right) &= \dfrac{\left(5+\sqrt{3}\right)\times\sqrt{3}}{\sqrt{3}\times\sqrt{3}}
  &&\justif{on multiplie en haut et en bas par $\sqrt{3}$}\\
h\!\left(2+\sqrt{3}\right) &= \res{\dfrac{5\sqrt{3}+3}{3}} &&\justif{$\sqrt{3}\times\sqrt{3}=3$}
\end{calculs}
\rem{Deux cas à reconnaître : un dénominateur qui est une \emph{somme} avec une racine se
multiplie par sa quantité conjuguée ; une racine \emph{seule} se multiplie par
elle-même. Erreur fréquente : ne multiplier que le dénominateur ; on multiplie
\emph{en haut et en bas}, sinon la valeur change.}
\end{exo}

\partie{Factoriser une expression}

\begin{exo}[Facteur commun]
\explique{On cherche ce qui se trouve dans \emph{chaque} terme (un nombre, $x$, ou les
deux), on le fait apparaître, puis on applique $ka+kb=k(a+b)$.}{méthode 4, exemple a)}
\begin{calculs}
A &= 5x+20 &&\justif{$20=5\times 4$ : $5$ est dans les deux termes}\\
A &= 5\times x+5\times 4 &&\\
A &= \res{5(x+4)} &&
\end{calculs}
\begin{calculs}
B &= 7x^{2}+3x &&\justif{$x^{2}=x\times x$ : $x$ est dans les deux termes}\\
B &= x\times 7x+x\times 3 &&\\
B &= \res{x(7x+3)} &&\justif{$7$ et $3$ n'ont pas de diviseur commun}
\end{calculs}
\begin{calculs}
C &= 18x^{2}-12x &&\justif{$6$ divise $18$ et $12$, et $x$ est dans les deux termes}\\
C &= 6x\times 3x-6x\times 2 &&\\
C &= \res{6x(3x-2)} &&
\end{calculs}
\begin{calculs}
D &= x^{2}+x &&\justif{le second terme $x$ s'écrit $x\times 1$}\\
D &= x\times x+x\times 1 &&\\
D &= \res{x(x+1)} &&
\end{calculs}
\begin{calculs}
E &= 4x^{2}-x &&\justif{même chose : $x=x\times 1$}\\
E &= x\times 4x-x\times 1 &&\\
E &= \res{x(4x-1)} &&
\end{calculs}
\begin{calculs}
F &= 2x+x(x+5) &&\justif{$2x=x\times 2$ : $x$ est dans les deux termes}\\
F &= x\times 2+x(x+5) &&\\
F &= x\bigl[2+(x+5)\bigr] &&\justif{on écrit ce qui reste de chaque terme}\\
F &= \res{x(x+7)} &&\justif{on réduit le crochet}
\end{calculs}
\rem{Erreurs fréquentes : en D et E, oublier le $1$ et écrire $x(x)$ ou $x(4x)$ ; en C,
s'arrêter à $2x(9x-6)$, juste mais incomplet ($9x-6$ contient encore le facteur $3$).
Contrôle : on redéveloppe et on doit retrouver l'expression de départ.}
\end{exo}

\begin{exo}[Identités remarquables et facteur commun]
\explique{Pour choisir l'outil : une parenthèse répétée, c'est un facteur commun ; deux
carrés séparés par un \og $-$\fg, c'est $a^{2}-b^{2}$ ; trois termes dont deux carrés,
c'est peut-être $a^{2}\pm 2ab+b^{2}$ (on vérifie le double produit).}{les identités
remarquables lues de droite à gauche}
\cas{A}
\begin{calculs}
A &= 9x^{2}+30x+25 &&\justif{pas de facteur commun ; $9x^{2}=(3x)^{2}$ et $25=5^{2}$}\\
A &= (3x)^{2}+2\times 3x\times 5+5^{2} &&\justif{double produit : $2\times 3x\times 5=30x$,
c'est bien le terme du milieu}\\
A &= \res{(3x+5)^{2}} &&\justif{$a^{2}+2ab+b^{2}=(a+b)^{2}$, avec $a=3x$ et $b=5$}
\end{calculs}
\cas{B}
\begin{calculs}
B &= 49x^{2}-36 &&\justif{deux carrés séparés par un \og $-$\fg}\\
B &= (7x)^{2}-6^{2} &&\justif{forme $a^{2}-b^{2}$, avec $a=7x$ et $b=6$}\\
B &= \res{(7x-6)(7x+6)} &&\justif{$(a-b)(a+b)$}
\end{calculs}
\cas{C}
\begin{calculs}
C &= (3x+1)(2x-5)+(3x+1)(x+7) &&\justif{facteur commun : la parenthèse entière $(3x+1)$}\\
C &= (3x+1)\bigl[(2x-5)+(x+7)\bigr] &&\justif{on écrit ce qui reste de chaque terme}\\
C &= \res{(3x+1)(3x+2)} &&\justif{on réduit : $2x+x=3x$ et $-5+7=2$}
\end{calculs}
\cas{D}
\begin{calculs}
D &= (4x-1)(x+6)-(4x-1)(3x-2) &&\justif{facteur commun : $(4x-1)$}\\
D &= (4x-1)\bigl[(x+6)-(3x-2)\bigr] &&\justif{des crochets : le \og $-$\fg{} porte sur
toute la seconde parenthèse}\\
D &= (4x-1)(x+6-3x+2) &&\justif{le \og $-$\fg{} change les signes de $3x-2$}\\
D &= \res{(4x-1)(-2x+8)} &&\justif{on réduit}
\end{calculs}
\rem{On peut aller plus loin : $-2x+8=2(4-x)$, d'où $D=2(4x-1)(4-x)$. Les deux réponses
sont justes.}
\cas{E}
\begin{calculs}
E &= (3x+2)^{2}-(x-4)^{2} &&\justif{forme $a^{2}-b^{2}$, avec $a=3x+2$ et $b=x-4$}\\
E &= \bigl[(3x+2)-(x-4)\bigr]\bigl[(3x+2)+(x-4)\bigr] &&\justif{$(a-b)(a+b)$, avec des
crochets}\\
E &= (3x+2-x+4)(3x+2+x-4) &&\justif{le \og $-$\fg{} change les signes de $x-4$}\\
E &= (2x+6)(4x-2) &&\justif{on réduit chaque facteur}\\
E &= 2(x+3)\times 2(2x-1) &&\justif{$2$ en facteur dans chaque parenthèse}\\
E &= \res{4(x+3)(2x-1)} &&
\end{calculs}
\cas{F}
\begin{calculs}
F &= 36-(x+2)^{2} &&\justif{$36=6^{2}$ : forme $a^{2}-b^{2}$, avec $a=6$ et $b=x+2$}\\
F &= \bigl[6-(x+2)\bigr]\bigl[6+(x+2)\bigr] &&\justif{$(a-b)(a+b)$}\\
F &= (6-x-2)(6+x+2) &&\justif{le \og $-$\fg{} change les signes de $x+2$}\\
F &= \res{(4-x)(x+8)} &&
\end{calculs}
\rem{Erreur fréquente en D, E et F : oublier que le \og $-$\fg{} change \emph{les deux}
signes de la parenthèse qui suit, par exemple écrire $3x+2-x-4$ en E. Ne pas
développer $(3x+2)^{2}$ et $(x-4)^{2}$ : on reconnaît directement $a^{2}-b^{2}$.}
\end{exo}

\begin{exo}[Faire apparaître un facteur commun]
\explique{Pour A, B et C, rien n'est commun à première vue. Mais chaque expression
contient une différence de carrés : en la factorisant, on voit apparaître un facteur qui
se trouve aussi dans l'autre terme.}{on factorise une partie, puis le tout}
\cas{A}
\begin{calculs}
A &= 9x^{2}-16-(3x+4)(x-2) &&\justif{$9x^{2}-16=(3x)^{2}-4^{2}$}\\
A &= (3x-4)(3x+4)-(3x+4)(x-2) &&\justif{$(a-b)(a+b)$ : le facteur $(3x+4)$ apparaît}\\
A &= (3x+4)\bigl[(3x-4)-(x-2)\bigr] &&\justif{facteur commun $(3x+4)$}\\
A &= (3x+4)(3x-4-x+2) &&\justif{le \og $-$\fg{} change les signes de $x-2$}\\
A &= (3x+4)(2x-2) &&\\
A &= \res{2(3x+4)(x-1)} &&\justif{$2$ en facteur dans $2x-2$}
\end{calculs}
\cas{B}
\begin{calculs}
B &= x^{2}-25-(x+5)(2x-3) &&\justif{$x^{2}-25=x^{2}-5^{2}$}\\
B &= (x-5)(x+5)-(x+5)(2x-3) &&\justif{le facteur $(x+5)$ apparaît}\\
B &= (x+5)\bigl[(x-5)-(2x-3)\bigr] &&\\
B &= (x+5)(x-5-2x+3) &&\\
B &= (x+5)(-x-2) &&\justif{on réduit}\\
B &= \res{-(x+5)(x+2)} &&\justif{facultatif : $-x-2=-(x+2)$}
\end{calculs}
\cas{C}
\begin{calculs}
C &= (x+3)(4x-1)-\left(x^{2}-9\right) &&\justif{$x^{2}-9=x^{2}-3^{2}$}\\
C &= (x+3)(4x-1)-(x-3)(x+3) &&\justif{le facteur $(x+3)$ apparaît}\\
C &= (x+3)\bigl[(4x-1)-(x-3)\bigr] &&\\
C &= (x+3)(4x-1-x+3) &&\\
C &= \res{(x+3)(3x+2)} &&
\end{calculs}
\cas{D}
\begin{calculs}
D &= x^{2}-7 &&\justif{$7$ n'est pas le carré d'un entier, mais
$7=\left(\sqrt{7}\right)^{2}$}\\
D &= x^{2}-\left(\sqrt{7}\right)^{2} &&\justif{forme $a^{2}-b^{2}$, avec $a=x$ et
$b=\sqrt{7}$}\\
D &= \res{\left(x-\sqrt{7}\right)\left(x+\sqrt{7}\right)} &&
\end{calculs}
\cas{E}
\begin{calculs}
E &= x^{2}+2\sqrt{5}\,x+5 &&\justif{trois termes : $x^{2}$ et $5=\left(\sqrt{5}\right)^{2}$
sont des carrés}\\
E &= x^{2}+2\times x\times\sqrt{5}+\left(\sqrt{5}\right)^{2} &&\justif{le double produit
$2\times x\times\sqrt{5}$ est bien le terme du milieu}\\
E &= \res{\left(x+\sqrt{5}\right)^{2}} &&\justif{$a^{2}+2ab+b^{2}$, avec $a=x$ et
$b=\sqrt{5}$}
\end{calculs}
\rem{À retenir : quand une expression contient une différence de carrés ($x^{2}-9$ ou $9x^{2}-16$) à côté d'une
parenthèse, on factorise la différence de carrés pour chercher cette parenthèse. En D
et E, $\sqrt{7}$ et $\sqrt{5}$ jouent le rôle de $b$ dans l'identité remarquable.}
\end{exo}

\partie{Équations}

\begin{exo}[Équations du premier degré]
\explique{On développe s'il y a des parenthèses, on regroupe les $x$ dans un membre et les
nombres dans l'autre, puis on divise par le coefficient de $x$.}{méthode 5, exemple a)}
\cas{a)}
\begin{calculs}
4x-9=3x+2 &\iff 4x-3x=2+9 &&\justif{les $x$ à gauche, les nombres à droite}\\
          &\iff x=11 &&
\end{calculs}
\explique{$S=\res{\ens{11}}$.}{vérification : $4\times 11-9=35$ et $3\times 11+2=35$}
\cas{b)}
\begin{calculs}
6x-5=x+20 &\iff 6x-x=20+5 &&\justif{les $x$ à gauche, les nombres à droite}\\
          &\iff 5x=25 &&\\
          &\iff x=5 &&\justif{on divise par $5$}
\end{calculs}
\explique{$S=\res{\ens{5}}$.}{}
\cas{c)}
\begin{calculs}
3x+4=3x-1 &\iff 3x-3x=-1-4 &&\\
          &\iff 0=-5 &&\justif{les $x$ ont disparu}
\end{calculs}
\explique{L'égalité $0=-5$ est fausse, quelle que soit la valeur de $x$ : l'équation n'a
pas de solution, $S=\res{\varnothing}$.}{il ne reste qu'une égalité entre deux
nombres}
\cas{d)}
\begin{calculs}
3(x-2)=2(-x+7) &\iff 3x-6=-2x+14 &&\justif{on développe chaque membre}\\
               &\iff 3x+2x=14+6 &&\justif{$-2x$ passe à gauche en $+2x$}\\
               &\iff 5x=20 &&\\
               &\iff x=4 &&
\end{calculs}
\explique{$S=\res{\ens{4}}$.}{vérification : $3\times 2=6$ et $2\times 3=6$}
\cas{e)}
\begin{calculs}
x-(2-5x)=4x+9 &\iff x-2+5x=4x+9 &&\justif{le \og $-$\fg{} change les signes de $2-5x$}\\
              &\iff 6x-2=4x+9 &&\\
              &\iff 6x-4x=9+2 &&\\
              &\iff 2x=11 &&\\
              &\iff x=\dfrac{11}{2} &&\justif{on divise par $2$ ; on garde la fraction}
\end{calculs}
\explique{$S=\res{\ens{\dfrac{11}{2}}}$.}{}
\cas{f)}
\begin{calculs}
4(x-1)+3(x+5)=0 &\iff 4x-4+3x+15=0 &&\justif{on développe}\\
                &\iff 7x+11=0 &&\justif{on réduit}\\
                &\iff 7x=-11 &&\\
                &\iff x=-\dfrac{11}{7} &&
\end{calculs}
\explique{$S=\res{\ens{-\dfrac{11}{7}}}$.}{valeur exacte : pas d'arrondi}
\rem{En c), il ne faut pas écrire $S=\ens{0}$ : $0$ n'est pas une solution
($3\times 0+4\ne 3\times 0-1$). Quand les $x$ disparaissent, on conclut selon l'égalité
qui reste : fausse, aucune solution ; vraie (comme $0=0$), tous les réels sont
solutions.}
\end{exo}

\begin{exo}[Produit nul et carrés]
\explique{Ici, on ne développe pas : un produit égal à $0$ se résout facteur par facteur
(propriété 2), et $x^{2}=k$ se résout directement (propriété 3).}{méthode 5,
exemples b) et c)}
\cas{a)}
\begin{calculs}
(x-6)(4+x)=0 &\iff x-6=0 \;\text{ ou }\; 4+x=0 &&\justif{produit nul : un des facteurs
est nul}\\
             &\iff x=6 \;\text{ ou }\; x=-4 &&
\end{calculs}
\explique{$S=\res{\ens{-4\,;6}}$.}{solutions rangées dans l'ordre croissant}
\cas{b)}
\begin{calculs}
5x(3x-7)=0 &\iff 5x=0 \;\text{ ou }\; 3x-7=0 &&\justif{deux facteurs : $5x$ et
$3x-7$}\\
           &\iff x=0 \;\text{ ou }\; x=\dfrac{7}{3} &&
\end{calculs}
\explique{$S=\res{\ens{0\,;\dfrac{7}{3}}}$.}{ne pas oublier la solution $0$}
\cas{c)}
\begin{calculs}
(2x+9)^{2}=0 &\iff 2x+9=0 &&\justif{$(2x+9)(2x+9)=0$ : les deux facteurs sont
identiques}\\
             &\iff x=-\dfrac{9}{2} &&
\end{calculs}
\explique{$S=\res{\ens{-\dfrac{9}{2}}}$.}{une seule solution}
\cas{d)}
\begin{calculs}
x^{2}=13 &\iff x=-\sqrt{13} \;\text{ ou }\; x=\sqrt{13} &&\justif{propriété 3 : $13>0$,
deux solutions}
\end{calculs}
\explique{$S=\res{\ens{-\sqrt{13}\,;\sqrt{13}}}$.}{on garde les valeurs exactes}
\cas{e)}
\begin{calculs}
3x^{2}-21=0 &\iff 3x^{2}=21 &&\justif{on isole $x^{2}$}\\
            &\iff x^{2}=7 &&\justif{on divise par $3$}\\
            &\iff x=-\sqrt{7} \;\text{ ou }\; x=\sqrt{7} &&\justif{propriété 3 : $7>0$}
\end{calculs}
\explique{$S=\res{\ens{-\sqrt{7}\,;\sqrt{7}}}$.}{}
\cas{f)}
\begin{calculs}
(3x+1)^{2}=25 &\iff 3x+1=5 \;\text{ ou }\; 3x+1=-5 &&\justif{propriété 3, avec $3x+1$ à
la place de $x$ : $\sqrt{25}=5$}\\
              &\iff 3x=4 \;\text{ ou }\; 3x=-6 &&\\
              &\iff x=\dfrac{4}{3} \;\text{ ou }\; x=-2 &&
\end{calculs}
\explique{$S=\res{\ens{-2\,;\dfrac{4}{3}}}$.}{vérification : $(-6+1)^{2}=25$ et
$(4+1)^{2}=25$}
\rem{Erreur fréquente en d), e) et f) : oublier la solution négative. En f), développer
$(3x+1)^{2}$ serait une fausse piste : on garde le carré et on applique la
propriété 3.}
\end{exo}

\begin{exo}[Mise en équation]
\cas{1)}
\explique{Soit $x$ le nombre de croissants vendus. Il y a deux fois plus de pains au
chocolat, soit $2x$, et $10$ chaussons de moins que de croissants, soit $x-10$.}{
\textbf{inconnue} : la quantité à partir de laquelle les autres s'expriment}
\begin{calculs}
x+2x+(x-10)=150 &\iff 4x-10=150 &&\justif{\textbf{équation} : le total des trois
sortes}\\
                &\iff 4x=160 &&\\
                &\iff x=40 &&
\end{calculs}
\explique{Pains au chocolat : $2\times 40=80$ ; chaussons : $40-10=30$ ; total :
$40+80+30=150$.}{\textbf{vérification} avec l'énoncé}
\explique{La boulangerie a vendu \res{$40$ croissants, $80$ pains au chocolat et $30$
chaussons aux pommes}.}{\textbf{conclusion} : une phrase qui répond à la question}
\cas{2)}
\explique{Soit $x$ le côté du carré, en cm. Le rectangle mesure $x+6$ sur $x-4$ : son
aire est $(x+6)(x-4)$, celle du carré est $x^{2}$.}{\textbf{inconnue}, puis traduction
des deux aires}
\begin{calculs}
(x+6)(x-4)=x^{2} &\iff x^{2}-4x+6x-24=x^{2} &&\justif{\textbf{équation} : les deux aires
sont égales ; on développe}\\
                 &\iff 2x-24=0 &&\justif{on retire $x^{2}$ des deux membres}\\
                 &\iff 2x=24 &&\\
                 &\iff x=12 &&
\end{calculs}
\explique{Le rectangle mesure $12+6=18$ cm sur $12-4=8$ cm. Son aire vaut
$18\times 8=144$ cm$^{2}$, comme celle du carré : $12^{2}=144$.}{\textbf{vérification} ;
la largeur $x-4$ est bien positive}
\explique{Le carré a pour côté \res{$12$ cm}, et le rectangle mesure \res{$18$ cm sur
$8$ cm}.}{\textbf{conclusion}}
\rem{En 2), l'équation contient $x^{2}$ dans les deux membres, mais ces termes
s'éliminent : elle se ramène au premier degré. Avant de conclure, on vérifie que la
solution a un sens (une longueur positive).}
\end{exo}

\partie{Fractions rationnelles}

\begin{exo}[Les quatre opérations]
\explique{Dans les quatre expressions, le seul dénominateur qui contient $x$ est $x$ : la
valeur interdite est \res{$0$}.}{$x=0$ annule le dénominateur}
\cas{A}
\begin{calculs}
A &= \dfrac{3}{4}+\dfrac{5}{x} &&\justif{addition : il faut un dénominateur commun,
$4x$}\\
A &= \dfrac{3\fois{x}}{4\fois{x}}+\dfrac{5\fois{4}}{x\fois{4}} &&\justif{chaque fraction
reçoit ce qui lui manque}\\
A &= \dfrac{3x}{4x}+\dfrac{20}{4x} &&\\
A &= \res{\dfrac{3x+20}{4x}} &&\justif{on additionne les numérateurs}
\end{calculs}
\cas{B}
\begin{calculs}
B &= \dfrac{3}{4}-\dfrac{5}{x} &&\justif{même dénominateur commun $4x$}\\
B &= \dfrac{3\fois{x}}{4\fois{x}}-\dfrac{5\fois{4}}{x\fois{4}} &&\\
B &= \res{\dfrac{3x-20}{4x}} &&
\end{calculs}
\cas{C}
\begin{calculs}
C &= \dfrac{3}{4}\times\dfrac{5}{x} &&\justif{produit : pas besoin de dénominateur
commun}\\
C &= \dfrac{3\times 5}{4\times x} &&\justif{numérateurs entre eux, dénominateurs entre
eux}\\
C &= \res{\dfrac{15}{4x}} &&
\end{calculs}
\cas{D}
\begin{calculs}
D &= \dfrac{3}{4}\div\dfrac{5}{x} &&\justif{diviser par $\frac{5}{x}$, c'est multiplier
par son inverse $\frac{x}{5}$}\\
D &= \dfrac{3}{4}\times\dfrac{x}{5} &&\\
D &= \res{\dfrac{3x}{20}} &&
\end{calculs}
\rem{Additionner ou soustraire : dénominateur commun obligatoire. Multiplier ou diviser :
on n'en a pas besoin. Erreur fréquente : $\frac{3}{4}+\frac{5}{x}=\frac{8}{4+x}$ (on
n'additionne jamais les dénominateurs). En D, on divise par $\frac{5}{x}$, qui doit être
non nul : c'est toujours le cas, car son numérateur $5$ n'est pas nul.}
\end{exo}

\begin{exo}[Valeurs interdites]
\explique{Valeurs interdites : on résout \og dénominateur $=0$\fg. Simplifier : on
factorise en haut et en bas, puis on barre un \emph{facteur} commun, jamais un
terme.}{méthode 7 du cours}
\cas{$A$}
\explique{$x-4=0$ pour $x=4$ : valeur interdite \res{$4$}. Le numérateur $x+5$ et le
dénominateur $x-4$ n'ont pas de facteur commun : \res{pas de simplification}.}{on ne
barre pas les $x$ : ce sont des termes d'une somme}
\cas{$B$}
\explique{$x(x-2)=0$ pour $x=0$ ou $x=2$ : valeurs interdites \res{$0$ et $2$}. Pas de
simplification : le numérateur $6$ ne contient ni $x$ ni $x-2$.}{produit nul :
propriété 2}
\cas{$C$}
\explique{$2x+5=0$ pour $x=-\dfrac{5}{2}$ : valeur interdite \res{$-\dfrac{5}{2}$}. Pas de
simplification.}{$2x=-5$, puis on divise par $2$}
\cas{$D$}
\explique{$x-4=0$ pour $x=4$ : valeur interdite \res{$4$}.}{on commence par les valeurs
interdites}
\begin{calculs}
D &= \dfrac{x^{2}-16}{x-4} &&\justif{pour $x\ne 4$}\\
D &= \dfrac{\barre{(x-4)}(x+4)}{\barre{(x-4)}} &&\justif{$x^{2}-16=x^{2}-4^{2}$ :
différence de carrés ; on barre le facteur commun}\\
D &= \res{x+4} &&
\end{calculs}
\explique{La valeur interdite reste $4$ : l'égalité $D=x+4$ n'est vraie que pour
$x\ne 4$.}{simplifier ne supprime pas une valeur interdite}
\cas{$E$}
\explique{$x^{2}+4=0$ revient à $x^{2}=-4$, qui n'a pas de solution : le dénominateur ne
s'annule jamais. \res{Aucune valeur interdite}, et pas de simplification.}{propriété 3 :
un carré n'est jamais négatif}
\cas{$F$}
\explique{$7-x=0$ pour $x=7$ : valeur interdite \res{$7$}. Pas de simplification.}{}
\rem{Erreur fréquente : simplifier $\frac{x+5}{x-4}$ en barrant les $x$. On ne barre que
des facteurs (des morceaux d'un produit), jamais des termes (des morceaux d'une
somme).}
\end{exo}

\begin{exo}[Réduire au même dénominateur]
\explique{On commence par les valeurs interdites ; puis chaque fraction reçoit, en haut et
en bas, ce qui manque à son dénominateur ; on développe le numérateur seulement, et on
garde le dénominateur factorisé.}{méthode 8 du cours}
\cas{$A$}
\explique{Valeur interdite : $-3$.}{$x+3=0$ pour $x=-3$}
\begin{calculs}
A &= 2+\dfrac{7}{x+3} &&\justif{on recopie}\\
A &= \dfrac{2\fois{(x+3)}}{x+3}+\dfrac{7}{x+3} &&\justif{même dénominateur :
$2=\frac{2(x+3)}{x+3}$}\\
A &= \dfrac{2x+6+7}{x+3} &&\justif{on développe le numérateur seulement}\\
A &= \res{\dfrac{2x+13}{x+3}} &&\justif{on réduit}
\end{calculs}
\cas{$B$}
\explique{Valeur interdite : $-1$.}{}
\begin{calculs}
B &= 4+\dfrac{x-2}{x+1} &&\\
B &= \dfrac{4\fois{(x+1)}}{x+1}+\dfrac{x-2}{x+1} &&\justif{$4=\frac{4(x+1)}{x+1}$}\\
B &= \dfrac{4x+4+x-2}{x+1} &&\\
B &= \res{\dfrac{5x+2}{x+1}} &&
\end{calculs}
\cas{$C$}
\explique{Valeurs interdites : $1$ et $-3$.}{dénominateur commun : le produit
$(x-1)(x+3)$}
\begin{calculs}
C &= \dfrac{x}{x-1}-\dfrac{2x}{x+3} &&\\
C &= \dfrac{x\fois{(x+3)}}{(x-1)\fois{(x+3)}}-\dfrac{2x\fois{(x-1)}}{(x+3)\fois{(x-1)}}
  &&\justif{chaque fraction reçoit ce qui lui manque}\\
C &= \dfrac{x(x+3)-2x(x-1)}{(x-1)(x+3)} &&\justif{le \og $-$\fg{} porte sur tout le
produit $2x(x-1)$}\\
C &= \dfrac{x^{2}+3x-2x^{2}+2x}{(x-1)(x+3)} &&\justif{$-2x(x-1)=-2x^{2}+2x$}\\
C &= \res{\dfrac{-x^{2}+5x}{(x-1)(x+3)}} &&\justif{on réduit le numérateur}
\end{calculs}
\cas{$D$}
\explique{Valeurs interdites : $-1$ et $\dfrac{3}{2}$.}{$2x-3=0$ pour $x=\frac{3}{2}$}
\begin{calculs}
D &= \dfrac{4}{x+1}-\dfrac{3}{2x-3} &&\\
D &= \dfrac{4\fois{(2x-3)}}{(x+1)\fois{(2x-3)}}-\dfrac{3\fois{(x+1)}}{(2x-3)\fois{(x+1)}} &&\\
D &= \dfrac{8x-12-3x-3}{(x+1)(2x-3)} &&\justif{$-3(x+1)=-3x-3$}\\
D &= \res{\dfrac{5x-15}{(x+1)(2x-3)}} &&
\end{calculs}
\cas{$E$}
\explique{Valeur interdite : $2$.}{}
\begin{calculs}
E &= 3x-\dfrac{2x+1}{x-2} &&\\
E &= \dfrac{3x\fois{(x-2)}}{x-2}-\dfrac{2x+1}{x-2} &&\justif{$3x=\frac{3x(x-2)}{x-2}$}\\
E &= \dfrac{3x^{2}-6x-(2x+1)}{x-2} &&\justif{on garde la parenthèse derrière le
\og $-$\fg}\\
E &= \dfrac{3x^{2}-6x-2x-1}{x-2} &&\justif{le \og $-$\fg{} change les signes de $2x+1$}\\
E &= \res{\dfrac{3x^{2}-8x-1}{x-2}} &&
\end{calculs}
\cas{$F$}
\explique{On remarque que $3x-12=3(x-4)$ : les deux dénominateurs ont le facteur $x-4$.
Valeur interdite : $4$ seulement ; dénominateur commun : $3(x-4)$.}{on factorise un
dénominateur avant de chercher le dénominateur commun}
\begin{calculs}
F &= \dfrac{2x}{x-4}-\dfrac{x+3}{3(x-4)} &&\justif{on écrit $3x-12$ sous forme
factorisée}\\
F &= \dfrac{2x\fois{3}}{(x-4)\fois{3}}-\dfrac{x+3}{3(x-4)} &&\justif{seule la première
fraction reçoit un facteur : $3$}\\
F &= \dfrac{6x-(x+3)}{3(x-4)} &&\\
F &= \dfrac{6x-x-3}{3(x-4)} &&\\
F &= \res{\dfrac{5x-3}{3(x-4)}} &&
\end{calculs}
\rem{En F, prendre $(x-4)(3x-12)$ comme dénominateur commun n'est pas faux, mais alourdit
tout : repérer $3x-12=3(x-4)$ fait gagner du temps. Erreur fréquente en C, E et F :
oublier que le \og $-$\fg{} devant une fraction porte sur \emph{tout} son numérateur.}
\end{exo}

\begin{exo}[Simplifier, multiplier, diviser]
\explique{Pour simplifier, on factorise le plus possible en haut et en bas ; un produit ou
un quotient de fractions se ramène à une seule fraction, que l'on simplifie
ensuite.}{on ne barre que des facteurs}
\cas{$A$}
\explique{$8(x+1)=0$ pour $x=-1$ : valeur interdite $-1$.}{}
\begin{calculs}
A &= \dfrac{4(x+1)^{2}}{8(x+1)} &&\justif{pour $x\ne -1$}\\
A &= \dfrac{\barre{4}\times\barre{(x+1)}\times(x+1)}{2\times\barre{4}\times\barre{(x+1)}}
  &&\justif{$(x+1)^{2}=(x+1)(x+1)$ et $8=2\times 4$}\\
A &= \res{\dfrac{x+1}{2}} &&
\end{calculs}
\cas{$B$}
\explique{Valeur interdite : $5$.}{}
\begin{calculs}
B &= \dfrac{25-x^{2}}{x-5} &&\justif{$25-x^{2}=5^{2}-x^{2}$ : différence de carrés}\\
B &= \dfrac{(5-x)(5+x)}{x-5} &&\justif{$5-x$ et $x-5$ sont opposés}\\
B &= \dfrac{-\barre{(x-5)}(x+5)}{\barre{(x-5)}} &&\justif{$5-x=-(x-5)$ : le facteur du
dénominateur apparaît}\\
B &= \res{-(x+5)} &&\justif{soit $-x-5$}
\end{calculs}
\cas{$C$}
\explique{Valeurs interdites : $2$ et $-2$.}{dénominateur commun $(x-2)(x+2)$}
\begin{calculs}
C &= \dfrac{x+2}{x-2}-\dfrac{x-2}{x+2} &&\\
C &= \dfrac{(x+2)\fois{(x+2)}}{(x-2)\fois{(x+2)}}-\dfrac{(x-2)\fois{(x-2)}}{(x+2)\fois{(x-2)}} &&\\
C &= \dfrac{(x+2)^{2}-(x-2)^{2}}{(x-2)(x+2)} &&\\
C &= \dfrac{x^{2}+4x+4-\left(x^{2}-4x+4\right)}{(x-2)(x+2)} &&\justif{deux identités
remarquables ; parenthèse derrière le \og $-$\fg}\\
C &= \dfrac{x^{2}+4x+4-x^{2}+4x-4}{(x-2)(x+2)} &&\\
C &= \res{\dfrac{8x}{(x-2)(x+2)}} &&
\end{calculs}
\rem{Autre voie pour le numérateur : $(x+2)^{2}-(x-2)^{2}$ est une différence de carrés,
égale à $\bigl[(x+2)-(x-2)\bigr]\bigl[(x+2)+(x-2)\bigr]$, soit $4\times 2x$ : on retrouve
$8x$.}
\cas{$D$}
\explique{Valeurs interdites : $-4$ et $-1$.}{les deux dénominateurs}
\begin{calculs}
D &= \dfrac{x^{2}-1}{x+4}\times\dfrac{x^{2}-16}{x+1} &&\justif{produit : numérateurs
entre eux, dénominateurs entre eux}\\
D &= \dfrac{(x-1)(x+1)(x-4)(x+4)}{(x+4)(x+1)} &&\justif{deux différences de carrés}\\
D &= \dfrac{(x-1)\barre{(x+1)}(x-4)\barre{(x+4)}}{\barre{(x+4)}\barre{(x+1)}}
  &&\justif{on barre les facteurs communs}\\
D &= \res{(x-1)(x-4)} &&
\end{calculs}
\cas{$E$}
\explique{$x^{2}+3x=x(x+3)$ s'annule en $0$ et en $-3$ ; $x+3$ s'annule en $-3$ ; enfin
on divise par $\dfrac{x+2}{x+3}$, qui doit être non nul : $x\ne -2$. Valeurs interdites :
$-3$ ; $-2$ et $0$.}{le diviseur ajoute une valeur interdite}
\begin{calculs}
E &= \dfrac{x^{2}-4}{x^{2}+3x}\div\dfrac{x+2}{x+3} &&\\
E &= \dfrac{x^{2}-4}{x^{2}+3x}\times\dfrac{x+3}{x+2} &&\justif{diviser, c'est
multiplier par l'inverse}\\
E &= \dfrac{(x-2)(x+2)(x+3)}{x(x+3)(x+2)} &&\justif{on factorise en haut et en bas}\\
E &= \dfrac{(x-2)\barre{(x+2)}\barre{(x+3)}}{x\barre{(x+3)}\barre{(x+2)}} &&\\
E &= \res{\dfrac{x-2}{x}} &&
\end{calculs}
\cas{$F$}
\explique{$x^{2}-9=(x-3)(x+3)$ : valeurs interdites \res{$-3$ et $3$}. Le numérateur
$2x+1$ n'a ni le facteur $x-3$ ni le facteur $x+3$ : \res{pas de
simplification}.}{on factorise pour comparer les facteurs}
\rem{En B, $5-x$ et $x-5$ sont \og à l'envers\fg : ils ne se simplifient pas tels quels, il
faut écrire $5-x=-(x-5)$. En E, la valeur $-2$ vient du diviseur : on ne divise pas par
une fraction nulle.}
\end{exo}

\partie{Équations quotient}

\begin{exo}[Quotient nul]
\explique{Propriété 4 : un quotient est nul si et seulement si son numérateur est nul, son
dénominateur ne l'étant pas. On cherche donc les valeurs interdites, puis on résout
\og numérateur $=0$\fg{} et on écarte les valeurs interdites.}{méthode 9, exemple a)}
\cas{a)}
\explique{Valeur interdite : $-2$.}{$x+2=0$ pour $x=-2$}
\begin{calculs}
4x-12=0 &\iff 4x=12 &&\justif{numérateur nul}\\
        &\iff x=3 &&
\end{calculs}
\explique{$3$ n'est pas interdit : $S=\res{\ens{3}}$.}{}
\cas{b)}
\explique{$5x^{2}=0$ pour $x=0$ : valeur interdite $0$.}{}
\begin{calculs}
3x+6=0 &\iff 3x=-6 &&\\
       &\iff x=-2 &&
\end{calculs}
\explique{$-2$ n'est pas interdit : $S=\res{\ens{-2}}$.}{}
\cas{c)}
\explique{Valeur interdite : $-1$.}{}
\begin{calculs}
(x-3)(2x+5)=0 &\iff x-3=0 \;\text{ ou }\; 2x+5=0 &&\justif{le numérateur est un produit
nul}\\
              &\iff x=3 \;\text{ ou }\; x=-\dfrac{5}{2} &&
\end{calculs}
\explique{Aucune des deux valeurs n'est interdite : $S=\res{\ens{-\dfrac{5}{2}\,;3}}$.}{}
\cas{d)}
\explique{$x^{2}\ge 0$, donc $x^{2}+3\ge 3$ : le dénominateur ne s'annule jamais. Aucune
valeur interdite.}{un carré n'est jamais négatif}
\begin{calculs}
7-2x=0 &\iff -2x=-7 &&\\
       &\iff x=\dfrac{7}{2} &&\justif{on divise par $-2$}
\end{calculs}
\explique{$S=\res{\ens{\dfrac{7}{2}}}$.}{}
\rem{Erreur fréquente : résoudre aussi \og dénominateur $=0$\fg{} et ajouter ces valeurs aux
solutions. Le dénominateur ne sert qu'à écarter des valeurs, jamais à en fournir.}
\end{exo}

\begin{exo}[Se ramener à un quotient nul]
\explique{Quand le second membre n'est pas $0$, on fait tout passer dans un membre, on
écrit un seul quotient (méthode 8), puis on applique la propriété 4.}{méthode 9,
exemple b)}
\cas{a)}
\explique{Valeur interdite : $-4$.}{}
\begin{calculs}
\dfrac{3x+2}{x+4}=2 &\iff \dfrac{3x+2}{x+4}-2=0 &&\justif{tout dans un membre}\\
                    &\iff \dfrac{3x+2-2(x+4)}{x+4}=0 &&\justif{un seul quotient :
$2=\frac{2(x+4)}{x+4}$}\\
                    &\iff \dfrac{3x+2-2x-8}{x+4}=0 &&\\
                    &\iff \dfrac{x-6}{x+4}=0 &&
\end{calculs}
\explique{$x-6=0$ pour $x=6$, qui n'est pas interdit : $S=\res{\ens{6}}$.}{numérateur nul,
valeur non interdite}
\cas{b)}
\explique{Valeur interdite : $-2$.}{}
\begin{calculs}
\dfrac{-3x}{x+2}=1 &\iff \dfrac{-3x}{x+2}-1=0 &&\justif{tout dans un membre}\\
                   &\iff \dfrac{-3x-(x+2)}{x+2}=0 &&\justif{$1=\frac{x+2}{x+2}$}\\
                   &\iff \dfrac{-4x-2}{x+2}=0 &&
\end{calculs}
\explique{$-4x-2=0$ pour $x=-\dfrac{1}{2}$, qui n'est pas interdit :
$S=\res{\ens{-\dfrac{1}{2}}}$.}{$-4x=2$, puis on divise par $-4$}
\cas{c)}
\explique{Valeur interdite : $0$.}{dénominateur commun : $3x$}
\begin{calculs}
\dfrac{1}{x}+\dfrac{1}{3}=1 &\iff \dfrac{1}{x}+\dfrac{1}{3}-1=0 &&\justif{tout dans un
membre}\\
  &\iff \dfrac{3}{3x}+\dfrac{x}{3x}-\dfrac{3x}{3x}=0 &&\justif{chaque terme sur $3x$}\\
  &\iff \dfrac{3-2x}{3x}=0 &&
\end{calculs}
\explique{$3-2x=0$ pour $x=\dfrac{3}{2}$, qui n'est pas interdit :
$S=\res{\ens{\dfrac{3}{2}}}$.}{vérification : $\frac{2}{3}+\frac{1}{3}=1$}
\cas{d)}
\explique{Valeur interdite : $6$.}{}
\begin{calculs}
\dfrac{5}{x-6}+1=0 &\iff \dfrac{5+(x-6)}{x-6}=0 &&\justif{déjà un membre nul ;
$1=\frac{x-6}{x-6}$}\\
                   &\iff \dfrac{x-1}{x-6}=0 &&
\end{calculs}
\explique{$x-1=0$ pour $x=1$, qui n'est pas interdit : $S=\res{\ens{1}}$.}{}
\cas{e)}
\explique{Valeur interdite : $-2$.}{}
\begin{calculs}
3-\dfrac{8-x}{x+2}=0 &\iff \dfrac{3(x+2)-(8-x)}{x+2}=0 &&\justif{on garde la parenthèse
derrière le \og $-$\fg}\\
                     &\iff \dfrac{3x+6-8+x}{x+2}=0 &&\\
                     &\iff \dfrac{4x-2}{x+2}=0 &&
\end{calculs}
\explique{$4x-2=0$ pour $x=\dfrac{1}{2}$, qui n'est pas interdit :
$S=\res{\ens{\dfrac{1}{2}}}$.}{}
\cas{f)}
\explique{Valeur interdite : $4$.}{le quotient est déjà égal à $0$}
\begin{calculs}
x^{2}-16=0 &\iff x^{2}=16 &&\justif{numérateur nul}\\
           &\iff x=-4 \;\text{ ou }\; x=4 &&\justif{propriété 3}
\end{calculs}
\explique{Mais $4$ est \emph{interdit} : on l'écarte. $S=\res{\ens{-4}}$.}{le quotient
n'existe pas pour $x=4$}
\rem{En f), une solution du numérateur est une valeur interdite : on l'écarte. C'est
pour cela qu'on écrit toujours les valeurs interdites en premier.}
\end{exo}

\begin{exo}[Égalité de deux quotients]
\explique{Méthode du cours : on fait tout passer dans un membre, on réduit au même
dénominateur, puis quotient nul. (Le produit en croix donnerait aussi le résultat ici,
mais il ne sert plus pour les inéquations de la partie suivante.)}{méthode 9,
exemple c)}
\cas{a)}
\explique{Valeurs interdites : $-2$ et $0$.}{dénominateur commun $x(x+2)$}
\begin{calculs}
\dfrac{5}{x+2}=\dfrac{3}{x} &\iff \dfrac{5}{x+2}-\dfrac{3}{x}=0 &&\justif{tout dans un
membre}\\
  &\iff \dfrac{5x-3(x+2)}{x(x+2)}=0 &&\justif{chaque fraction reçoit ce qui lui manque}\\
  &\iff \dfrac{2x-6}{x(x+2)}=0 &&\justif{$5x-3x-6=2x-6$}
\end{calculs}
\explique{$2x-6=0$ pour $x=3$, qui n'est pas interdit : $S=\res{\ens{3}}$.}{vérification :
$\frac{5}{5}=1$ et $\frac{3}{3}=1$}
\cas{b)}
\explique{Valeurs interdites : $0$ et $-1$.}{}
\begin{calculs}
\dfrac{2x+1}{x}=\dfrac{2x+5}{x+1} &\iff \dfrac{2x+1}{x}-\dfrac{2x+5}{x+1}=0 &&\\
  &\iff \dfrac{(2x+1)(x+1)-x(2x+5)}{x(x+1)}=0 &&
\end{calculs}
\explique{On calcule le numérateur à part.}{calcul intermédiaire}
\begin{calculs}
(2x+1)(x+1)-x(2x+5) &= 2x^{2}+2x+x+1-2x^{2}-5x &&\justif{$-x(2x+5)=-2x^{2}-5x$}\\
(2x+1)(x+1)-x(2x+5) &= -2x+1 &&\justif{les $x^{2}$ s'éliminent}
\end{calculs}
\explique{L'équation équivaut à $\dfrac{-2x+1}{x(x+1)}=0$ ; $-2x+1=0$ pour
$x=\dfrac{1}{2}$, qui n'est pas interdit : $S=\res{\ens{\dfrac{1}{2}}}$.}{}
\cas{c)}
\explique{Valeurs interdites : $-1$ et $2$.}{}
\begin{calculs}
\dfrac{3x-2}{x+1}=\dfrac{3x+1}{x-2} &\iff \dfrac{3x-2}{x+1}-\dfrac{3x+1}{x-2}=0 &&\\
  &\iff \dfrac{(3x-2)(x-2)-(3x+1)(x+1)}{(x+1)(x-2)}=0 &&
\end{calculs}
\explique{On calcule le numérateur à part.}{}
\begin{calculs}
(3x-2)(x-2)-(3x+1)(x+1) &= \left(3x^{2}-8x+4\right)-\left(3x^{2}+4x+1\right) &&\justif{chaque
produit entre parenthèses}\\
(3x-2)(x-2)-(3x+1)(x+1) &= 3x^{2}-8x+4-3x^{2}-4x-1 &&\\
(3x-2)(x-2)-(3x+1)(x+1) &= -12x+3 &&
\end{calculs}
\explique{$-12x+3=0$ pour $x=\dfrac{3}{12}$, soit $x=\dfrac{1}{4}$, qui n'est pas
interdit : $S=\res{\ens{\dfrac{1}{4}}}$.}{on simplifie la fraction par $3$}
\cas{d)}
\explique{Valeur interdite : $2$.}{un seul dénominateur, $x-2$}
\begin{calculs}
\dfrac{x+4}{x-2}-3=\dfrac{8-x}{x-2}
  &\iff \dfrac{x+4}{x-2}-3-\dfrac{8-x}{x-2}=0 &&\justif{tout dans un membre}\\
  &\iff \dfrac{x+4-3(x-2)-(8-x)}{x-2}=0 &&\justif{$3=\frac{3(x-2)}{x-2}$}\\
  &\iff \dfrac{x+4-3x+6-8+x}{x-2}=0 &&\\
  &\iff \dfrac{-x+2}{x-2}=0 &&
\end{calculs}
\explique{$-x+2=0$ pour $x=2$, qui est \emph{interdit} : l'équation n'a pas de solution,
$S=\res{\varnothing}$.}{la seule valeur trouvée est écartée}
\cas{e)}
\explique{Valeur interdite : $1$.}{les deux dénominateurs sont égaux}
\begin{calculs}
\dfrac{x+5}{x-1}=\dfrac{3}{x-1} &\iff \dfrac{x+5-3}{x-1}=0 &&\justif{même dénominateur :
on soustrait les numérateurs}\\
  &\iff \dfrac{x+2}{x-1}=0 &&
\end{calculs}
\explique{$x+2=0$ pour $x=-2$, qui n'est pas interdit : $S=\res{\ens{-2}}$.}{}
\cas{f)}
\explique{Valeurs interdites : $1$ et $3$. Calcul long : on réduit d'abord le membre de
gauche.}{$x-1=0$ pour $x=1$ ; $3-x=0$ pour $x=3$}
\begin{calculs}
1-\dfrac{x+1}{x-1} &= \dfrac{x-1}{x-1}-\dfrac{x+1}{x-1} &&\justif{$1=\frac{x-1}{x-1}$}\\
1-\dfrac{x+1}{x-1} &= \dfrac{x-1-x-1}{x-1} &&\justif{le \og $-$\fg{} change les signes de
$x+1$}\\
1-\dfrac{x+1}{x-1} &= \dfrac{-2}{x-1} &&
\end{calculs}
\explique{Puis on résout l'équation obtenue.}{dénominateur commun $(x-1)(3-x)$}
\begin{calculs}
\dfrac{-2}{x-1}=\dfrac{4}{3-x}
  &\iff \dfrac{-2}{x-1}-\dfrac{4}{3-x}=0 &&\\
  &\iff \dfrac{-2(3-x)-4(x-1)}{(x-1)(3-x)}=0 &&\\
  &\iff \dfrac{-6+2x-4x+4}{(x-1)(3-x)}=0 &&\\
  &\iff \dfrac{-2x-2}{(x-1)(3-x)}=0 &&
\end{calculs}
\explique{$-2x-2=0$ pour $x=-1$, qui n'est pas interdit : $S=\res{\ens{-1}}$.}{vérification :
$1-\frac{0}{-2}=1$ et $\frac{4}{4}=1$}
\rem{En d), la seule valeur trouvée est interdite : pas de solution. En e), les
dénominateurs sont égaux : on compare directement les numérateurs, en gardant la valeur
interdite.}
\end{exo}

\partie{Inéquations et tableaux de signes}

\begin{exo}[Inéquations du premier degré]
\explique{On isole $x$ comme dans une équation. Une seule règle de plus : multiplier ou
diviser par un nombre \emph{négatif} change le sens de l'inégalité.}{méthode 10,
exemple a)}
\cas{a)}
\begin{calculs}
3x+12>0 &\iff 3x>-12 &&\\
        &\iff x>-4 &&\justif{on divise par $3>0$ : le sens ne change pas}
\end{calculs}
\explique{$S=\res{\left]-4\,;+\infty\right[}$.}{$-4$ exclu : inégalité stricte}
\cas{b)}
\begin{calculs}
4x-1<9-x &\iff 4x+x<9+1 &&\justif{les $x$ à gauche, les nombres à droite}\\
         &\iff 5x<10 &&\\
         &\iff x<2 &&\justif{on divise par $5>0$}
\end{calculs}
\explique{$S=\res{\left]-\infty\,;2\right[}$.}{}
\cas{c)}
\begin{calculs}
3(x-1)\ge 7-2x &\iff 3x-3\ge 7-2x &&\justif{on développe}\\
               &\iff 5x\ge 10 &&\\
               &\iff x\ge 2 &&
\end{calculs}
\explique{$S=\res{\left[2\,;+\infty\right[}$.}{$2$ compris : inégalité large}
\cas{d)}
\begin{calculs}
\dfrac{2-5x}{3}>0 &\iff 2-5x>0 &&\justif{on multiplie par $3>0$ : le sens ne change pas}\\
                  &\iff -5x>-2 &&\\
                  &\iff x<\dfrac{2}{5} &&\justif{on divise par $-5<0$ : le sens change}
\end{calculs}
\explique{$S=\res{\left]-\infty\,;\dfrac{2}{5}\right[}$.}{}
\cas{e)}
\begin{calculs}
3(x+2)<3x+4 &\iff 3x+6<3x+4 &&\\
            &\iff 6<4 &&\justif{les $x$ disparaissent}
\end{calculs}
\explique{L'inégalité $6<4$ est fausse, quel que soit $x$ : $S=\res{\varnothing}$.}{comme
en c) de l'exercice 10}
\cas{f)}
\begin{calculs}
\dfrac{x-3}{2}-\dfrac{1-x}{4}>0 &\iff 2(x-3)-(1-x)>0 &&\justif{on multiplie par $4>0$ :
$4\times\frac{x-3}{2}=2(x-3)$}\\
  &\iff 2x-6-1+x>0 &&\justif{le \og $-$\fg{} change les signes de $1-x$}\\
  &\iff 3x-7>0 &&\\
  &\iff x>\dfrac{7}{3} &&
\end{calculs}
\explique{$S=\res{\left]\dfrac{7}{3}\,;+\infty\right[}$.}{}
\rem{Seules la multiplication et la division par un nombre \emph{négatif} changent le
sens. Erreur fréquente en d) : garder $x>\frac{2}{5}$. Contrôle : $x=0$ vérifie
$\frac{2-0}{3}>0$, donc $0$ doit être dans $S$.}
\end{exo}

\begin{exo}[Tableaux de signes]
\explique{Une ligne par facteur $ax+b$ : il s'annule en $-\dfrac{b}{a}$, a le signe de $a$
après cette valeur et le signe contraire avant. La dernière ligne s'obtient par la règle
des signes, colonne par colonne.}{propriété 5 du cours}
\cas{A}
\explique{$x+1$ s'annule en $-1$, $x-5$ en $5$ ; coefficients de $x$ positifs : chaque
facteur est négatif avant sa valeur, positif après.}{valeurs rangées dans l'ordre
croissant}
\begin{center}
\begin{tikzpicture}[scale=0.85]
  \tkzTabInit[lgt=2.6,espcl=2]{$x$/1,$x+1$/1,$x-5$/1,$(x-5)(x+1)$/1.2}{$-\infty$,$-1$,$5$,$+\infty$}
  \tkzTabLine{,-,z,+,t,+,}
  \tkzTabLine{,-,t,-,z,+,}
  \tkzTabLine{,+,z,-,z,+,}
\end{tikzpicture}
\end{center}
\cas{B}
\explique{$2x+1$ s'annule en $-\dfrac{1}{2}$, avec $a=2>0$. $3-x$ s'annule en $3$, avec
$a=-1<0$ : positif avant $3$, négatif après.}{attention au facteur $3-x$ : le coefficient
de $x$ est négatif}
\begin{center}
\begin{tikzpicture}[scale=0.85]
  \tkzTabInit[lgt=2.6,espcl=2]{$x$/1,$2x+1$/1,$3-x$/1,$(3-x)(2x+1)$/1.2}{$-\infty$,$-\frac{1}{2}$,$3$,$+\infty$}
  \tkzTabLine{,-,z,+,t,+,}
  \tkzTabLine{,+,t,+,z,-,}
  \tkzTabLine{,-,z,+,z,-,}
\end{tikzpicture}
\end{center}
\cas{C}
\explique{On factorise d'abord : $x^{2}-16=(x-4)(x+4)$. Les facteurs s'annulent en $-4$ et
en $4$.}{un tableau de signes se fait sur un produit de facteurs $ax+b$}
\begin{center}
\begin{tikzpicture}[scale=0.85]
  \tkzTabInit[lgt=2.6,espcl=2]{$x$/1,$x+4$/1,$x-4$/1,$x^{2}-16$/1.2}{$-\infty$,$-4$,$4$,$+\infty$}
  \tkzTabLine{,-,z,+,t,+,}
  \tkzTabLine{,-,t,-,z,+,}
  \tkzTabLine{,+,z,-,z,+,}
\end{tikzpicture}
\end{center}
\cas{D}
\explique{Trois facteurs : $4x+1$ s'annule en $-\dfrac{1}{4}$, $2x$ en $0$, $x-3$ en $3$ ;
tous ont un coefficient positif.}{à gauche, trois signes $-$ : le produit est négatif}
\begin{center}
\begin{tikzpicture}[scale=0.85]
  \tkzTabInit[lgt=3,espcl=1.9]{$x$/1,$4x+1$/1,$2x$/1,$x-3$/1,$2x(x-3)(4x+1)$/1.2}{$-\infty$,$-\frac{1}{4}$,$0$,$3$,$+\infty$}
  \tkzTabLine{,-,z,+,t,+,t,+,}
  \tkzTabLine{,-,t,-,z,+,t,+,}
  \tkzTabLine{,-,t,-,t,-,z,+,}
  \tkzTabLine{,-,z,+,z,-,z,+,}
\end{tikzpicture}
\end{center}
\cas{E}
\explique{Valeur interdite : $1$ (double barre). $x+4$ s'annule en $-4$, avec $a=1>0$ ;
$1-x$ s'annule en $1$, avec $a=-1<0$.}{la règle des signes vaut pour un quotient comme
pour un produit}
\begin{center}
\begin{tikzpicture}[scale=0.85]
  \tkzTabInit[lgt=2.6,espcl=2]{$x$/1,$x+4$/1,$1-x$/1,$\frac{x+4}{1-x}$/1.4}{$-\infty$,$-4$,$1$,$+\infty$}
  \tkzTabLine{,-,z,+,t,+,}
  \tkzTabLine{,+,t,+,z,-,}
  \tkzTabLine{,-,z,+,d,-,}
\end{tikzpicture}
\end{center}
\cas{F}
\explique{Valeur interdite : $-\dfrac{3}{2}$. Au numérateur, $x$ s'annule en $0$ et
$x-2$ en $2$ ; trois coefficients positifs.}{$2x+3=0$ pour $x=-\frac{3}{2}$}
\begin{center}
\begin{tikzpicture}[scale=0.85]
  \tkzTabInit[lgt=3,espcl=1.9]{$x$/1,$2x+3$/1,$x$/1,$x-2$/1,$\frac{x(x-2)}{2x+3}$/1.4}{$-\infty$,$-\frac{3}{2}$,$0$,$2$,$+\infty$}
  \tkzTabLine{,-,z,+,t,+,t,+,}
  \tkzTabLine{,-,t,-,z,+,t,+,}
  \tkzTabLine{,-,t,-,t,-,z,+,}
  \tkzTabLine{,-,d,+,z,-,z,+,}
\end{tikzpicture}
\end{center}
\rem{Sous une valeur interdite, une double barre, pas un zéro. Contrôle de la dernière
ligne : chaque fois qu'on franchit une valeur, un seul facteur change de signe, donc le
produit (ou le quotient) change de signe.}
\end{exo}

\begin{exo}[Inéquations produit et quotient]
\explique{On compare un produit ou un quotient à $0$ : on dresse son tableau de signes,
puis on lit les intervalles. Une inégalité large garde les zéros du numérateur ; une
valeur interdite est toujours exclue.}{méthode 10, exemples b) et c)}
\cas{a)}
\explique{Pas de dénominateur : aucune valeur interdite. On factorise
$9x^{2}-4=(3x-2)(3x+2)$ pour n'avoir que des facteurs $ax+b$ ; ils s'annulent en
$-\dfrac{2}{3}$, $\dfrac{2}{3}$ et $2$.}{$9x^{2}-4=(3x)^{2}-2^{2}$}
\begin{center}
\begin{tikzpicture}[scale=0.85]
  \tkzTabInit[lgt=3.8,espcl=1.9]{$x$/1,$3x+2$/1,$3x-2$/1,$x-2$/1,$\left(9x^{2}-4\right)(x-2)$/1.2}{$-\infty$,$-\frac{2}{3}$,$\frac{2}{3}$,$2$,$+\infty$}
  \tkzTabLine{,-,z,+,t,+,t,+,}
  \tkzTabLine{,-,t,-,z,+,t,+,}
  \tkzTabLine{,-,t,-,t,-,z,+,}
  \tkzTabLine{,-,z,+,z,-,z,+,}
\end{tikzpicture}
\end{center}
\explique{$S=\res{\left[-\dfrac{2}{3}\,;\dfrac{2}{3}\right]\cup\left[2\,;+\infty\right[}$.}{on
lit les signes $+$ et les zéros ($\ge$ : bornes comprises)}
\cas{b)}
\explique{$x^{2}-4=(x-2)(x+2)$ : valeurs interdites $-2$ et $2$. Le numérateur $5-2x$
s'annule en $\dfrac{5}{2}$, avec $a=-2<0$.}{$5-2x$ : positif avant $\frac{5}{2}$, négatif
après}
\begin{center}
\begin{tikzpicture}[scale=0.85]
  \tkzTabInit[lgt=3,espcl=1.9]{$x$/1,$5-2x$/1,$x+2$/1,$x-2$/1,$\frac{5-2x}{x^{2}-4}$/1.4}{$-\infty$,$-2$,$2$,$\frac{5}{2}$,$+\infty$}
  \tkzTabLine{,+,t,+,t,+,z,-,}
  \tkzTabLine{,-,z,+,t,+,t,+,}
  \tkzTabLine{,-,t,-,z,+,t,+,}
  \tkzTabLine{,+,d,-,d,+,z,-,}
\end{tikzpicture}
\end{center}
\explique{$S=\res{\left]-2\,;2\right[\cup\left[\dfrac{5}{2}\,;+\infty\right[}$.}{signes $-$
et zéro : $\frac{5}{2}$ compris, $-2$ et $2$ exclus (interdits)}
\cas{c)}
\explique{Valeur interdite : $-\dfrac{1}{3}$. Au numérateur, $x^{2}-9=(x-3)(x+3)$ s'annule
en $-3$ et en $3$.}{$3x+1=0$ pour $x=-\frac{1}{3}$}
\begin{center}
\begin{tikzpicture}[scale=0.85]
  \tkzTabInit[lgt=3,espcl=1.9]{$x$/1,$x+3$/1,$3x+1$/1,$x-3$/1,$\frac{x^{2}-9}{3x+1}$/1.4}{$-\infty$,$-3$,$-\frac{1}{3}$,$3$,$+\infty$}
  \tkzTabLine{,-,z,+,t,+,t,+,}
  \tkzTabLine{,-,t,-,z,+,t,+,}
  \tkzTabLine{,-,t,-,t,-,z,+,}
  \tkzTabLine{,-,z,+,d,-,z,+,}
\end{tikzpicture}
\end{center}
\explique{$S=\res{\left[-3\,;-\dfrac{1}{3}\right[\cup\left[3\,;+\infty\right[}$.}{$-3$ et $3$
compris ; $-\frac{1}{3}$ exclu (interdit)}
\cas{d)}
\explique{Valeur interdite : $2$. Le second membre n'est pas $0$ : on s'y ramène, comme
pour une équation quotient.}{méthode 10, exemple c)}
\begin{calculs}
\dfrac{x+3}{2-x}\le 1 &\iff \dfrac{x+3}{2-x}-1\le 0 &&\justif{tout dans un membre}\\
                      &\iff \dfrac{x+3-(2-x)}{2-x}\le 0 &&\justif{un seul quotient
comparé à $0$}\\
                      &\iff \dfrac{2x+1}{2-x}\le 0 &&
\end{calculs}
\explique{$2x+1$ s'annule en $-\dfrac{1}{2}$ ($a=2>0$) ; $2-x$ s'annule en $2$
($a=-1<0$).}{propriété 5}
\begin{center}
\begin{tikzpicture}[scale=0.85]
  \tkzTabInit[lgt=2.6,espcl=2]{$x$/1,$2x+1$/1,$2-x$/1,$\frac{2x+1}{2-x}$/1.4}{$-\infty$,$-\frac{1}{2}$,$2$,$+\infty$}
  \tkzTabLine{,-,z,+,t,+,}
  \tkzTabLine{,+,t,+,z,-,}
  \tkzTabLine{,-,z,+,d,-,}
\end{tikzpicture}
\end{center}
\explique{$S=\res{\left]-\infty\,;-\dfrac{1}{2}\right]\cup\left]2\,;+\infty\right[}$.}{signes
$-$ et zéro ; $2$ exclu (interdit)}
\cas{e)}
\explique{Valeur interdite : $-4$.}{}
\begin{calculs}
\dfrac{2-x}{x+4}>3 &\iff \dfrac{2-x}{x+4}-3>0 &&\justif{tout dans un membre}\\
                   &\iff \dfrac{2-x-3(x+4)}{x+4}>0 &&\\
                   &\iff \dfrac{2-x-3x-12}{x+4}>0 &&\\
                   &\iff \dfrac{-4x-10}{x+4}>0 &&
\end{calculs}
\explique{$-4x-10$ s'annule en $-\dfrac{5}{2}$ et $a=-4<0$ : positif avant
$-\dfrac{5}{2}$, négatif après.}{$-4x=10$, donc $x=-\frac{10}{4}$}
\begin{center}
\begin{tikzpicture}[scale=0.85]
  \tkzTabInit[lgt=2.6,espcl=2]{$x$/1,$x+4$/1,$-4x-10$/1,$\frac{-4x-10}{x+4}$/1.4}{$-\infty$,$-4$,$-\frac{5}{2}$,$+\infty$}
  \tkzTabLine{,-,z,+,t,+,}
  \tkzTabLine{,+,t,+,z,-,}
  \tkzTabLine{,-,d,+,z,-,}
\end{tikzpicture}
\end{center}
\explique{$S=\res{\left]-4\,;-\dfrac{5}{2}\right[}$.}{signes $+$ seulement : inégalité
stricte}
\cas{f)}
\explique{Valeurs interdites : $-2$ et $-3$.}{}
\begin{calculs}
\dfrac{x+1}{x+2}\le\dfrac{x-1}{x+3}
  &\iff \dfrac{x+1}{x+2}-\dfrac{x-1}{x+3}\le 0 &&\justif{tout dans un membre}\\
  &\iff \dfrac{(x+1)(x+3)-(x-1)(x+2)}{(x+2)(x+3)}\le 0 &&
\end{calculs}
\explique{On calcule le numérateur à part.}{calcul intermédiaire}
\begin{calculs}
(x+1)(x+3)-(x-1)(x+2) &= \left(x^{2}+4x+3\right)-\left(x^{2}+x-2\right) &&\\
(x+1)(x+3)-(x-1)(x+2) &= x^{2}+4x+3-x^{2}-x+2 &&\\
(x+1)(x+3)-(x-1)(x+2) &= 3x+5 &&
\end{calculs}
\explique{On étudie donc $\dfrac{3x+5}{(x+2)(x+3)}\le 0$ : $3x+5$ s'annule en
$-\dfrac{5}{3}$, $x+2$ en $-2$, $x+3$ en $-3$.}{ordre : $-3<-2<-\frac{5}{3}$}
\begin{center}
\begin{tikzpicture}[scale=0.85]
  \tkzTabInit[lgt=3,espcl=1.9]{$x$/1,$x+3$/1,$x+2$/1,$3x+5$/1,$\frac{3x+5}{(x+2)(x+3)}$/1.4}{$-\infty$,$-3$,$-2$,$-\frac{5}{3}$,$+\infty$}
  \tkzTabLine{,-,z,+,t,+,t,+,}
  \tkzTabLine{,-,t,-,z,+,t,+,}
  \tkzTabLine{,-,t,-,t,-,z,+,}
  \tkzTabLine{,-,d,+,d,-,z,+,}
\end{tikzpicture}
\end{center}
\explique{$S=\res{\left]-\infty\,;-3\right[\cup\left]-2\,;-\dfrac{5}{3}\right]}$.}{$-\frac{5}{3}$
compris ; $-3$ et $-2$ exclus}
\rem{Une valeur interdite est toujours exclue de $S$, même avec $\le$ ou $\ge$. Erreur
fréquente en d), e) et f) : multiplier les deux membres par le dénominateur. Son signe
dépend de $x$, on ne sait donc pas si le sens change : c'est pour cela qu'on se ramène
à un quotient comparé à $0$.}
\end{exo}

\partie{Composition de polynômes}

\begin{exo}[Premières compositions]
\explique{$(P\circ Q)(x)=P\bigl(Q(x)\bigr)$ : on recopie $P(x)$ en remplaçant chaque $x$
par $Q(x)$, \emph{entre parenthèses}, puis on développe et on réduit.}{méthode 11 du
cours}
\cas{a)}
\begin{calculs}
(P\circ Q)(x) &= 3(2x+1)+2 &&\justif{on recopie $P(x)=3x+2$ : $x$ devient $(2x+1)$}\\
(P\circ Q)(x) &= 6x+3+2 &&\justif{on distribue le $3$}\\
(P\circ Q)(x) &= \res{6x+5} &&
\end{calculs}
\explique{Contrôle avec $1$ : $Q(1)=3$ et $P(3)=11$ ; on trouve aussi $6+5=11$.}{$(P\circ
Q)(1)$ doit être égal à $P\bigl(Q(1)\bigr)$}
\cas{b)}
\begin{calculs}
(P\circ Q)(x) &= 2(3x+2)+1 &&\justif{on recopie $P(x)=2x+1$ : $x$ devient $(3x+2)$}\\
(P\circ Q)(x) &= 6x+4+1 &&\\
(P\circ Q)(x) &= \res{6x+5} &&
\end{calculs}
\cas{c)}
\begin{calculs}
(P\circ Q)(x) &= 3\left(2x^{2}+x-4\right)-2 &&\justif{$x$ devient $\left(2x^{2}+x-4\right)$ :
la parenthèse est indispensable}\\
(P\circ Q)(x) &= 6x^{2}+3x-12-2 &&\justif{le $3$ multiplie les trois termes}\\
(P\circ Q)(x) &= \res{6x^{2}+3x-14} &&
\end{calculs}
\cas{Comparaison}
\explique{En b), on a échangé les rôles des deux fonctions de a) : on a calculé, avec les
fonctions de a), $Q\circ P$ au lieu de $P\circ Q$. On obtient le même polynôme $6x+5$ :
\res{pour ces deux fonctions, l'ordre ne change rien}.}{les fonctions de b) sont celles
de a), échangées}
\rem{C'est une coïncidence propre à ces deux fonctions : en général
$P\circ Q\ne Q\circ P$ (exemple du cours et exercice 25).}
\end{exo}

\begin{exo}[Composer et calculer]
\cas{1)}
\begin{calculs}
(P\circ Q)(x) &= (3x-1)^{2}+2(3x-1)-1 &&\justif{chaque $x$ de $P$ devient $(3x-1)$, y
compris dans $x^{2}$}\\
(P\circ Q)(x) &= 9x^{2}-6x+1+6x-2-1 &&\justif{$(a-b)^{2}$ avec $a=3x$ et $b=1$ ; on
distribue le $2$}\\
(P\circ Q)(x) &= \res{9x^{2}-2} &&\justif{les termes en $x$ s'annulent}
\end{calculs}
\cas{2)}
\begin{calculs}
(P\circ Q)(x) &= \left(x^{2}+1\right)^{2}-\left(x^{2}+1\right) &&\justif{chaque $x$ de $P$
devient $\left(x^{2}+1\right)$}\\
(P\circ Q)(x) &= x^{4}+2x^{2}+1-x^{2}-1 &&\justif{$(a+b)^{2}$ avec $a=x^{2}$ :
$\left(x^{2}\right)^{2}=x^{4}$ ; le \og $-$\fg{} change deux signes}\\
(P\circ Q)(x) &= \res{x^{4}+x^{2}} &&
\end{calculs}
\cas{3) a)}
\begin{calculs}
(P\circ Q)(x) &= (x+4)^{2}-3 &&\justif{$x$ devient $(x+4)$ dans $x^{2}-3$}\\
(P\circ Q)(x) &= x^{2}+8x+16-3 &&\justif{$(a+b)^{2}$}\\
(P\circ Q)(x) &= \res{x^{2}+8x+13} &&\\[4pt]
(P\circ Q)(-2) &= (-2)^{2}+8\times(-2)+13 &&\justif{on remplace $x$ par $-2$ dans le
résultat}\\
(P\circ Q)(-2) &= 4-16+13 &&\\
(P\circ Q)(-2) &= \res{1} &&
\end{calculs}
\cas{3) b)}
\begin{calculs}
Q(-2) &= -2+4 &&\justif{on calcule d'abord l'image par $Q$}\\
Q(-2) &= 2 &&\\
P(2) &= 2^{2}-3 &&\justif{puis l'image de ce résultat par $P$}\\
P(2) &= \res{1} &&
\end{calculs}
\explique{On retrouve la même valeur : $(P\circ Q)(-2)=P\bigl(Q(-2)\bigr)$.}{c'est la
définition de la composée}
\rem{Pour une seule valeur, la seconde voie est souvent plus rapide : on n'a pas besoin
de l'expression de $P\circ Q$.}
\end{exo}

\begin{exo}[Composer dans les deux sens]
\cas{1) a)}
\begin{calculs}
(P\circ Q)(x) &= \left(x^{2}-2x\right)^{2}+1 &&\justif{chaque $x$ de $P$ devient
$\left(x^{2}-2x\right)$}\\
(P\circ Q)(x) &= \left(x^{2}\right)^{2}-2\times x^{2}\times 2x+(2x)^{2}+1
  &&\justif{$(a-b)^{2}$, avec $a=x^{2}$ et $b=2x$}\\
(P\circ Q)(x) &= \res{x^{4}-4x^{3}+4x^{2}+1} &&\justif{$x^{2}\times x=x^{3}$}\\[4pt]
(Q\circ P)(x) &= \left(x^{2}+1\right)^{2}-2\left(x^{2}+1\right) &&\justif{cette fois,
chaque $x$ de $Q$ devient $\left(x^{2}+1\right)$}\\
(Q\circ P)(x) &= x^{4}+2x^{2}+1-2x^{2}-2 &&\justif{$(a+b)^{2}$ ; on distribue $-2$}\\
(Q\circ P)(x) &= \res{x^{4}-1} &&
\end{calculs}
\explique{Les deux composées sont différentes.}{l'ordre compte}
\cas{1) b)}
\explique{$P$ et $Q$ sont de degré \res{$2$} : leur plus haute puissance de $x$ est
$x^{2}$. $P\circ Q$ est de degré \res{$4$} : son terme de plus haut degré est
$x^{4}$, de coefficient $1\ne 0$.}{définition 5 du cours}
\rem{Le degré monte : en remplaçant $x$ par une expression de degré $2$ dans $x^{2}$, le
carré donne $x^{4}$.}
\cas{2)}
\begin{calculs}
(P\circ Q)(x) &= 2\left(\sqrt{2}\,x+1\right)^{2}-\left(\sqrt{2}\,x+1\right) &&\justif{chaque
$x$ de $P$ devient $\left(\sqrt{2}\,x+1\right)$}\\
(P\circ Q)(x) &= 2\left(2x^{2}+2\sqrt{2}\,x+1\right)-\sqrt{2}\,x-1
  &&\justif{$(a+b)^{2}$, avec $\left(\sqrt{2}\,x\right)^{2}=2x^{2}$}\\
(P\circ Q)(x) &= 4x^{2}+4\sqrt{2}\,x+2-\sqrt{2}\,x-1 &&\justif{on distribue le $2$}\\
(P\circ Q)(x) &= \res{4x^{2}+3\sqrt{2}\,x+1} &&\justif{$4\sqrt{2}\,x-\sqrt{2}\,x=3\sqrt{2}\,x$}
\end{calculs}
\rem{Erreur fréquente : $\left(\sqrt{2}\,x\right)^{2}=\sqrt{2}\,x^{2}$. Le carré porte sur
les deux facteurs : $\left(\sqrt{2}\right)^{2}\times x^{2}=2x^{2}$.}
\end{exo}

\end{document}
